% =============================================================================
% Rapport de projet Graphes et OpenData
% Analyse du Reseau WETH sur Polygon lors du Lundi Noir du 5 aout 2024
% Universite Paris Nanterre - Licence MIAGE
% =============================================================================

\documentclass[12pt,a4paper]{article}

% --- Encodage et langue ---
\usepackage[utf8]{inputenc}          % Encodage UTF-8 des sources
\usepackage[T1]{fontenc}             % Encodage T1 pour les caracteres francais
\usepackage[french]{babel}           % Conventions typographiques francaises
\usepackage{csquotes}                % Requis par biblatex avec babel

% --- Mathematiques ---
\usepackage{amsmath, amssymb}        % Notation mathematique (Q, G, C, alpha...)

% --- Figures et tableaux ---
\usepackage{graphicx}                % Inclusion de figures (\includegraphics)
\usepackage{float}                   % Placement [H] des flottants
\usepackage{booktabs}                % Tableaux professionnels (\toprule, \midrule, \bottomrule)
\usepackage{caption, subcaption}     % Personnalisation des legendes + sous-figures

% --- Couleurs et code ---
\usepackage{xcolor}                  % Definitions de couleurs
\usepackage{listings}                % Listings de code Python

% --- Mise en page ---
\usepackage{fancyhdr}                % En-tetes et pieds de page
\usepackage{titlesec}                % Formatage des titres de sections
\usepackage{tocloft}                 % Formatage de la table des matieres
\usepackage{setspace}                % Interligne

% --- Geometrie (marges conformes au template Paris Nanterre) ---
\usepackage[a4paper, left=2.5cm, right=2.5cm, top=2.5cm, bottom=2.5cm]{geometry}

% --- Bibliographie ---
\usepackage[backend=biber, style=numeric, sorting=none, defernumbers=true]{biblatex}

% --- References croisees (charger en dernier) ---
\usepackage[colorlinks=true, linkcolor=red, citecolor=green, urlcolor=cyan]{hyperref}
\usepackage[french]{cleveref}        % References intelligentes (apres hyperref)

% =============================================================================
% Configuration
% =============================================================================

% Chemin des figures (depuis rapport/ vers figures/ a la racine)
\graphicspath{{../figures/}}

% Fichier de bibliographie
\addbibresource{references.bib}

% --- Palette de couleurs du projet ---
\definecolor{primary}{HTML}{2563EB}
\definecolor{secondary}{HTML}{7C3AED}
\definecolor{accent}{HTML}{F59E0B}

% --- Configuration des listings Python ---
\lstset{
  language=Python,
  basicstyle=\ttfamily\small,
  keywordstyle=\color{primary}\bfseries,
  commentstyle=\color{secondary}\itshape,
  stringstyle=\color{accent},
  numbers=left,
  numberstyle=\tiny\color{gray},
  frame=single,
  breaklines=true,
  captionpos=b,
  tabsize=4,
  showstringspaces=false
}

% =============================================================================
% Document
% =============================================================================

\begin{document}

% ---- Page de garde ----
\begin{titlepage}
\thispagestyle{empty}

\begin{flushleft}
% Logo Paris Nanterre (placeholder si le fichier n'existe pas)
% Placer le fichier logo_nanterre.png dans le repertoire rapport/ pour l'afficher
\IfFileExists{logo_nanterre.png}{%
  \includegraphics[width=7cm]{logo_nanterre.png}%
}{%
  \textbf{[Logo Universit\'e Paris Nanterre]}%
}
\end{flushleft}

\vspace{1.5cm}

\begin{center}

{\large Licence MIAGE}

\vspace{1cm}

{\Large \textbf{Rapport de projet Graphes et OpenData}}

\vspace{1.5cm}

{\LARGE \textcolor{primary}{\textbf{Analyse du R\'eseau WETH sur Polygon\\lors du Lundi Noir du 5 ao\^ut 2024}}}

\vspace{1.5cm}

{\large Projet r\'ealis\'e du XX au XX}

\vfill

{\large \textbf{Membres du groupe}}

\vspace{0.5cm}

{\large
Nom Pr\'enom 1\\
Nom Pr\'enom 2\\
Nom Pr\'enom 3\\
}

\vspace{1cm}

{\large Ann\'ee universitaire 2024--2025}

\end{center}

\end{titlepage}

% ---- Table des matieres ----
\newpage
\tableofcontents

% ---- Corps du rapport ----
\newpage
\section*{Chapitre 14: Prototype Interactif - crypto-graph-explorer}
\addcontentsline{toc}{section}{Chapitre 14: Prototype Interactif}

Un prototype interactif a été développé pour explorer visuellement le réseau WETH sur Polygon.
Ce chapitre présente la mise en œuvre du prototype et ses fonctionnalités d'analyse.

\subsection*{14.1 Architecture et Technologies}

Le prototype est construit avec les technologies suivantes:
\begin{itemize}
  \item \textbf{Frontend}: React 18 + TypeScript avec Vite
  \item \textbf{Visualisation}: D3.js v7 pour le graphe force-directed
  \item \textbf{Physique}: Simulation multi-force (charge, lien, collision, centre)
  \item \textbf{Effets visuels}: Filtres SVG (glow, ombre, animations D3)
\end{itemize}

\subsection*{14.2 Analyses Implémentées}

Le prototype inclut 11 analyses différentes pour explorer le réseau:

\subsubsection*{Analyses Avancées (8 analyses)}
\begin{enumerate}
  \item \textbf{Clustering}: Détection de communautés avec classification par couleur
  \item \textbf{Détection de Ponts}: Identification des nœuds critiques (articulation points)
  \item \textbf{Centralité Betweenness}: Mesure de l'importance des nœuds intermédiaires
  \item \textbf{Distribution des Degrés}: Analyse statistique de la connectivité
  \item \textbf{Chemins Optimaux}: Calcul des plus courts chemins entre nœuds
  \item \textbf{Analyse de Similarité}: Mesure de similarité structurelle entre nœuds
  \item \textbf{K-Core}: Décomposition en noyaux denses (degré de cohésion)
  \item \textbf{Simulation de Flux}: Simulation de propagation via le réseau
\end{enumerate}

\subsubsection*{Questions de Recherche (4 RQ + 1 supplémentaire)}
\begin{enumerate}
  \item \textbf{RQ1 - Topologie}: Structure générale du réseau WETH
  \item \textbf{RQ2 - Centralité}: Identification des nœuds centraux et influents
  \item \textbf{RQ3 - Clustering}: Communautés et structures cohérentes\footnote{Implémenté mais intégré aux analyses avancées}
  \item \textbf{RQ4 - Temporalité}: Dynamiques temporelles (placeholder pour données temporelles)
  \item \textbf{RQ5 - Small-World}: Propriétés de petit monde et chemins courts
  \item \textbf{RQ6 - Flux Pondérés}: Analyse des flux avec poids PageRank
\end{enumerate}

\subsection*{14.3 Interaction Utilisateur}

Le prototype offre des interactions intuitives:
\begin{itemize}
  \item \textbf{Sélection de nœuds}: Clic pour isoler et analyser un nœud spécifique
  \item \textbf{Drag-and-drop}: Déplacement interactif de nœuds dans le graphe
  \item \textbf{Hover}: Affichage dynamique des connexions et informations
  \item \textbf{Zoom}: Navigation fluide du graphe avec zoom/pan
  \item \textbf{Suppression progressive}: Simulation d'impact de la suppression de nœuds
\end{itemize}

\subsection*{14.4 Visualisation et Esthétique}

Le design du prototype reprend les principes d'Epistemenet:
\begin{itemize}
  \item \textbf{Codage couleur}: Les nœuds sont colorés par communauté
  \item \textbf{Taille variable}: Proportionnelle au PageRank (importance)
  \item \textbf{Effets de glow}: Soulignement des sélections et connexions
  \item \textbf{Animations fluides}: Transitions et pulsations pour offrir un retour visuel
  \item \textbf{Statistiques en temps réel}: Panneau graphique des métriques du réseau
\end{itemize}

\clearpage
% sections/introduction.tex
% Phase 9: Introduction (RAPP-01)

\section{Introduction}
\label{sec:introduction}

Le 5 ao\^ut 2024, les march\'es financiers mondiaux ont c\'ed\'e en quelques
heures, un \'episode que la presse a appel\'e le \og Lundi Noir \fg{}. Le
d\'eclencheur est mon\'etaire~: le 31 juillet, la Banque du Japon rel\`eve ses
taux directeurs par surprise~\cite{boj2024rate}, ce qui force le d\'enouement
des positions de \textit{carry trade} sur le yen. Le m\'ecanisme est connu~:
des investisseurs empruntent en yen \`a taux quasi nul (le Japon maintient des
taux n\'egatifs depuis 2016) pour placer ailleurs \`a rendement plus
\'elev\'e~\cite{reuters2024carrytrade}. Quand le yen s'appr\'ecie brutalement,
tout le monde doit d\'eboucler en m\^eme temps pour couvrir les pertes de
change. Les ventes forc\'ees se r\'epercutent sur toutes les classes d'actifs,
y compris la crypto~: sur Polygon, blockchain de couche~2 d'Ethereum, les
protocoles de finance d\'ecentralis\'ee (DeFi) encaissent des vagues de
liquidations et de transferts de panique. Sur la seule journ\'ee du 5 ao\^ut,
plus de 401\,000 transferts du jeton WETH (Wrapped Ether) sont
enregistr\'es sur Polygon.

Ces transferts forment un r\'eseau orient\'e pond\'er\'e. Chaque adresse
Ethereum est un sommet, chaque transfert une ar\^ete orient\'ee de l'\'emetteur
vers le r\'ecepteur, et le montant transf\'er\'e d\'efinit le poids de l'ar\^ete.
La th\'eorie des graphes donne des outils pour analyser ce r\'eseau~:
topologie, acteurs centraux, structure communautaire, dynamique temporelle
des \'echanges~\cite{newman2018networks}. C'est plus utile qu'un simple
d\'ecompte de transactions, parce qu'on peut alors voir o\`u se concentrent
les flux, quels protocoles servent de relais, et comment le r\'eseau se
fragmente quand le march\'e tangue.

Le rapport s'organise autour de six questions de recherche (RQ)~:

\begin{enumerate}
    \item \textbf{RQ1 --- Topologie du r\'eseau~:} analyser la structure globale
    \`a travers la distribution des degr\'es, la densit\'e, les composantes connexes
    et le diam\`etre du graphe.
    \item \textbf{RQ2 --- Centralit\'e~:} identifier les n\oe{}uds les plus
    influents par quatre mesures de centralit\'e (degr\'e, \textit{betweenness},
    \textit{closeness}, PageRank) et interpr\'eter leurs r\^oles.
    \item \textbf{RQ3 --- Communaut\'es~:} d\'etecter la structure communautaire
    par l'algorithme de Louvain, mesurer la modularit\'e et analyser la
    distribution des tailles de communaut\'es.
    \item \textbf{RQ4 --- Dynamique temporelle~:} d\'ecouper la journ\'ee en
    fen\^etres horaires et suivre l'\'evolution des m\'etriques du r\'eseau pour
    identifier les pics d'activit\'e.
    \item \textbf{RQ5 --- Propri\'et\'es petit-monde~:} \'evaluer le coefficient
    de clustering et la longueur moyenne des chemins, puis comparer au mod\`ele
    al\'eatoire d'Erd\H{o}s-R\'enyi.
    \item \textbf{RQ6 --- Flux pond\'er\'es~:} analyser la distribution des
    volumes WETH \'echang\'es, quantifier la concentration par le coefficient de
    Gini et la courbe de Lorenz.
\end{enumerate}

L'environnement de travail et le jeu de donn\'ees sont d\'ecrits en premier,
puis les six questions sont trait\'ees une par une (m\'ethodologie,
r\'esultats, interpr\'etation). Le bilan revient sur ce que les six analyses
disent ensemble, et propose quelques pistes pour aller plus loin.

En compl\'ement, un prototype web ouvre les m\^emes donn\'ees \`a
l'exploration directe \`a la souris~: \url{https://crypto-graph-explorer.vercel.app}~\cite{cryptographexplorer_web}.
On peut y zoomer dans le graphe, filtrer par communaut\'e, ouvrir un
n\oe{}ud, et relancer les six RQ et les cinq analyses compl\'ementaires
(section~\ref{sec:complementaires}) sur des m\'etriques pr\'e-calcul\'ees.
Les chiffres exacts rapport\'es ici viennent de l'ex\'ecution Python
d\'ecrite dans l'annexe~\ref{ann:manuel}.

\clearpage
% sections/environnement.tex
% Phase 9: Environnement de travail (RAPP-02)

\section{Environnement de travail}
\label{sec:environnement}

Voici l'environnement mat\'eriel et logiciel utilis\'e pour le projet.

\begin{table}[H]
\centering
\small
\caption{Environnement de d\'eveloppement.}
\label{tab:environnement}
\begin{tabular}{p{4cm}lp{6cm}}
\toprule
\textbf{Outil} & \textbf{Version} & \textbf{R\^ole} \\
\midrule
Syst\`eme d'exploitation & Windows~11 & Plateforme de d\'eveloppement \\
Python & 3.10+ & Langage principal d'analyse \\
pip & 24.x & Gestionnaire de paquets \\
Git & 2.x & Contr\^ole de version \\
VS Code & 1.9x & \'Editeur de code \\
\TeX{} Live & 2024 & Distribution \LaTeX{} \\
pdf\LaTeX{} & (inclus) & Compilation du rapport \\
Biber & (inclus) & Backend bibliographique \\
latexmk & (inclus) & Automatisation compilation \\
Node.js & 20.x & Ex\'ecution du prototype web \\
npm & 10.x & Gestionnaire de paquets JavaScript \\
Vercel CLI & 51.x & D\'eploiement du prototype web \\
\bottomrule
\end{tabular}
\end{table}

Python a \'et\'e choisi pour ses biblioth\`eques scientifiques (NumPy, Pandas, NetworkX, Matplotlib). Les d\'ependances s'installent via \texttt{pip install -r requirements.txt}.

Le rapport est compil\'e avec pdf\LaTeX{}, plus simple que Lua\LaTeX{} ou Xe\LaTeX{} et suffisant pour un document en fran\c{c}ais avec l'encodage T1. \TeX{} Live fournit tous les paquets n\'ecessaires, et \texttt{latexmk} automatise les passes de compilation.

Le prototype web~\cite{cryptographexplorer_web} tourne sur Node.js et est b\^ati avec Vite. En local~: \texttt{npm install} puis \texttt{npm run dev}. Pour la mise en ligne~: \texttt{vercel deploy}. Les biblioth\`eques utilis\'ees c\^ot\'e front sont d\'ecrites dans la section~\ref{sec:bibliotheques}.

\clearpage
\input{sections/description}
\clearpage
% sections/bibliotheques.tex
% Phase 9: Bibliotheques, Outils et Technologies (RAPP-04)

\section{Biblioth\`eques, Outils et technologies}
\label{sec:bibliotheques}

Les biblioth\`eques et technologies utilis\'ees sont d\'ecrites ci-dessous,
organis\'ees par cat\'egorie.

\subsection{Analyse de graphe}

\paragraph{NetworkX (v3.4.2).}
NetworkX~\cite{networkx_docs} fournit les structures \texttt{DiGraph} et
\texttt{Graph} utilis\'ees dans ce projet, ainsi que les algorithmes de
centralit\'e (degr\'e, betweenness, closeness, PageRank), de d\'etection de
communaut\'es (Louvain, int\'egr\'e depuis la version~3.2), de composantes
connexes, de clustering et de plus courts chemins. C'est le socle commun
aux six questions de recherche.

\subsection{Visualisation}

\paragraph{Matplotlib (v3.9.4).}
Matplotlib~\cite{matplotlib_docs} g\'en\`ere les 16~figures PNG \`a 200~DPI.
Deux contextes graphiques sont utilis\'es~: fond sombre (\texttt{\#0F172A})
pour les graphes de r\'eseau, fond clair pour les graphiques analytiques.
La palette de couleurs (bleu \texttt{\#2563EB}, violet \texttt{\#7C3AED},
orange \texttt{\#F59E0B}) est centralis\'ee dans \texttt{src/style.py}.

\subsection{Traitement de donn\'ees}

\paragraph{Pandas (v2.2.3).}
Pandas charge et manipule le CSV de 401\,709 transactions~: conversion
des timestamps Unix en \texttt{datetime}, agr\'egation des ar\^etes multiples
via \texttt{groupby}, et fen\^etrage horaire pour l'analyse temporelle (RQ4).

\paragraph{NumPy (v2.1).}
NumPy fournit les tableaux et op\'erations vectoris\'ees pour les calculs
num\'eriques~: distributions de degr\'es, courbes de Lorenz, statistiques
descriptives (moyennes, m\'edianes, percentiles).

\paragraph{SciPy (v1.14).}
SciPy compl\`ete NumPy pour les tests statistiques, les comparaisons de
distributions et les corr\'elations de rang (Spearman).

\paragraph{powerlaw (v2.0.0).}
La biblioth\`eque \texttt{powerlaw} ajuste des distributions \`a queue
lourde~: estimation de l'exposant $\alpha$ et du seuil $x_{\min}$ par
maximum de vraisemblance, et tests de comparaison entre mod\`eles (loi de
puissance, lognormale, exponentielle), selon la m\'ethodologie de Clauset
et al.~\cite{clauset2009powerlaw}. Utilis\'ee en RQ1 et RQ6.

\subsection{Visualisation interactive (prototype web)}

\paragraph{D3.js (v7.9).}
D3.js calcule le \textit{layout} force-directed (\texttt{d3-force}) qui
positionne les $12\,880$ n\oe{}uds du prototype web~\cite{cryptographexplorer_web},
et g\`ere le zoom et le drag de la souris. On a pr\'ef\'er\'e D3 \`a Plotly
pour garder la main sur le rendu SVG et ne pas d\'ependre d'un serveur de
\textit{plotting}.

\paragraph{React (v18.3) et TypeScript (v5.8).}
React d\'ecoupe l'interface en composants (canevas du graphe, panneau
d'analyse, panneau de statistiques). TypeScript ajoute le typage des
structures de graphe et des m\'etriques par n\oe{}ud, ce qui \'evite des
incoh\'erences entre le CSV charg\'e et la simulation D3.

\paragraph{Vite (v5.4), TailwindCSS et shadcn/ui.}
Vite sert le code en d\'eveloppement et le compile pour la mise en ligne.
TailwindCSS donne les styles. La biblioth\`eque shadcn/ui (au-dessus de
Radix~UI) fournit les boutons, sliders et autres contr\^oles du panneau
d'analyse.

\subsection{R\'edaction et compilation}

\paragraph{\LaTeX{} avec pdf\LaTeX{}.}
Le rapport est compos\'e avec le syst\`eme \LaTeX{}, compil\'e par le moteur
pdf\LaTeX{}. Le document utilise la classe \texttt{article} avec 18~paquets
organis\'es selon un ordre de chargement strict~: \texttt{babel} pour les
conventions typographiques fran\c{c}aises, \texttt{amsmath} pour la notation
math\'ematique, \texttt{booktabs} pour les tableaux professionnels,
\texttt{graphicx} pour l'inclusion des figures, \texttt{biblatex} pour la
gestion bibliographique, \texttt{hyperref} pour les liens cliquables et
\texttt{cleveref} pour les r\'ef\'erences crois\'ees intelligentes. La
bibliographie est g\'er\'ee par \texttt{biblatex} avec le backend Biber,
offrant un support Unicode complet pour les noms d'auteurs comportant des
caract\`eres accentu\'es.

\clearpage
% sections/travail_realise.tex
% Phase 10: Travail realise (RAPP-05)

\section{Travail r\'ealis\'e}
\label{sec:travail_realise}

Le projet s'est d\'eroul\'e en cinq phases~: collecte des donn\'ees,
construction du graphe, analyses par question de recherche, r\'edaction
du rapport, et d\'eveloppement d'un prototype web interactif compl\'ementaire.

\paragraph{Collecte et pr\'eparation des donn\'ees.}
Les 401\,709 transferts WETH du 5 ao\^ut 2024 ont \'et\'e extraits de la
table publique \texttt{crypto\_polygon.transactions} sur Google
BigQuery~\cite{bigquery_public}. Le r\'esultat, stock\'e dans le fichier
CSV \texttt{polygon\_weth\_crash\_2024-08-05.csv} (74\,Mo), a \'et\'e soumis
\`a un pr\'etraitement d'agr\'egation~: les ar\^etes multiples entre une m\^eme
paire d'adresses (source, target) sont fusionn\'ees par somme des poids
(volume WETH) et comptage des transactions. Un filtrage
\textit{dust} \`a $10^{-12}$~WETH supprime les micro-transferts n\'egligeables.
Le graphe r\'esultant comprend 12\,880 sommets et 39\,999 ar\^etes orient\'ees
pond\'er\'ees.

\paragraph{Construction du graphe.}
Le module \texttt{graph\_builder.py} sert d'interface entre les donn\'ees
brutes et les six analyses. Sa fonction \texttt{load\_and\_build()} retourne
un triplet \texttt{(df, digraph, graph)}~: le DataFrame Pandas, un graphe
orient\'e pond\'er\'e (\texttt{DiGraph}), et un graphe non orient\'e
(\texttt{Graph}) pour RQ3 et RQ5. Toutes les analyses travaillent ainsi sur
le m\^eme jeu de donn\'ees.

\paragraph{Analyses par question de recherche.}
Six modules ind\'ependants (\texttt{src/rq1\_*.py} \`a \texttt{src/rq6\_*.py})
r\'ealisent chacun une analyse et produisent deux figures PNG \`a 200~DPI
(16~au total). Le module \texttt{src/style.py} assure la coh\'erence
graphique~: fond sombre pour les graphes de r\'eseau, fond clair pour les
graphiques analytiques. Le script \texttt{run\_all.py} ex\'ecute les six
analyses \`a la suite.

\paragraph{R\'edaction du rapport.}
Le rapport est r\'edig\'e en \LaTeX{} et compil\'e avec pdf\LaTeX{} + Biber.
Le fichier principal \texttt{main.tex} inclut 17~fichiers \texttt{.tex}
(sections et annexes). La bibliographie est scind\'ee en deux~: r\'ef\'erences
acad\'emiques et webographie, filtr\'ees par mot-cl\'e via \texttt{biblatex}.

\paragraph{Prototype web interactif.}
En parall\`ele du rapport, on a construit un prototype web pour explorer
le r\'eseau sans installer Python. Il est d\'eploy\'e \`a l'adresse
\url{https://crypto-graph-explorer.vercel.app}~\cite{cryptographexplorer_web}.
La visualisation force-directed utilise D3.js, l'interface est faite en
React + Vite avec les composants shadcn/ui. Le prototype lit un CSV de
m\'etriques pr\'e-calcul\'ees (\texttt{nodes\_metrics.csv}, $12\,880$
n\oe{}uds, degr\'e, PageRank et num\'ero de communaut\'e), et un panneau
lat\'eral donne acc\`es aux six RQ et aux cinq analyses compl\'ementaires
de la section~\ref{sec:complementaires}. Les chiffres exacts ($G$,
$\sigma$, $Q$) viennent uniquement de l'ex\'ecution Python sur les
donn\'ees brutes~; le prototype, lui, calcule des versions approch\'ees
adapt\'ees au navigateur. L'annexe~\ref{ann:manuel} d\'ecrit l'acc\`es
et ses limites plus en d\'etail.

\begin{table}[H]
\centering
\caption{R\'epartition des t\^aches du projet}
\label{tab:repartition}
\begin{tabular}{lll}
\toprule
\textbf{Membre} & \textbf{T\^ache} & \textbf{Contribution} \\
\midrule
Thibaud THOMAS-LAMOTTE & Collecte BigQuery, RQ1, RQ2, RQ3 & Topologie, centralit\'e, communaut\'es \\
Youcef HEMISSI & RQ4, RQ5, RQ6, rapport \LaTeX{} & Temporel, petit-monde, flux \\
\bottomrule
\end{tabular}
\end{table}

\clearpage
% sections/rq1.tex
% Phase 3: RQ1 Topologie (RQ1-01 a RQ1-06)

\section{RQ1 -- Analyse de la topologie du r\'eseau}
\label{sec:rq1}

On commence par examiner la structure topologique globale du r\'eseau~: distribution des degr\'es, densit\'e, composantes connexes, diam\`etre et longueur moyenne des chemins. Ces m\'etriques donnent une premi\`ere image de l'architecture des \'echanges WETH pendant le crash.

\subsection{M\'ethodologie}
\label{sec:rq1_methodo}

L'analyse topologique repose sur la construction d'un graphe orient\'e pond\'er\'e $G = (V, E)$ \`a partir des 401\,000 transferts WETH enregistr\'es le 5 ao\^ut 2024. Chaque adresse Ethereum constitue un sommet, et chaque transfert une ar\^ete orient\'ee du portefeuille \'emetteur vers le portefeuille r\'ecepteur. Les ar\^etes multiples entre une m\^eme paire de sommets sont agr\'eg\'ees par somme des poids (volume WETH) et comptage des transactions.

Quatre familles de m\'etriques sont calcul\'ees. Premi\`erement, la \textbf{densit\'e} du r\'eseau, d\'efinie pour un graphe orient\'e comme~:
\begin{equation}
d = \frac{|E|}{|V| \cdot (|V| - 1)}
\label{eq:density}
\end{equation}
o\`u $|V|$ d\'esigne le nombre de sommets et $|E|$ le nombre d'ar\^etes. Cette m\'etrique quantifie la proportion de connexions existantes par rapport au nombre maximal de connexions possibles. Le \textbf{degr\'e moyen} est donn\'e par~:
\begin{equation}
\bar{k} = \frac{|E|}{|V|}
\label{eq:avg_degree}
\end{equation}

Deuxi\`emement, les \textbf{distributions des degr\'es} entrants ($k_{\text{in}}$) et sortants ($k_{\text{out}}$) sont analys\'ees s\'epar\'ement. Plut\^ot qu'un histogramme classique, nous utilisons la fonction de r\'epartition compl\'ementaire cumulative (CCDF), d\'efinie par $P(X \geq k)$, repr\'esent\'ee en \'echelle logarithmique sur les deux axes. Cette repr\'esentation est pr\'ef\'erable pour les distributions \`a queue lourde car elle \'evite les artefacts li\'es au choix des intervalles de regroupement et r\'ev\`ele plus clairement la forme de la distribution dans les hautes valeurs de degr\'e~\cite{clauset2009powerlaw}. L'observation d'une relation lin\'eaire en \'echelle log-log sugg\`ere un comportement h\'et\'erog\`ene caract\'eristique des r\'eseaux dits \og sans \'echelle \fg{}, o\`u quelques n\oe{}uds concentrent un nombre disproportionn\'e de connexions~\cite{barabasi1999scalefree}.

Troisi\`emement, les \textbf{composantes connexes} sont identifi\'ees selon deux d\'efinitions. Les composantes faiblement connexes (WCC, \textit{weakly connected components}) ignorent l'orientation des ar\^etes~: deux sommets appartiennent \`a la m\^eme WCC s'il existe un chemin non orient\'e entre eux. Les composantes fortement connexes (SCC, \textit{strongly connected components}) respectent l'orientation~: deux sommets appartiennent \`a la m\^eme SCC si et seulement s'il existe un chemin orient\'e de l'un \`a l'autre et r\'eciproquement~\cite{newman2018networks}. La taille de la composante g\'eante, exprim\'ee en pourcentage du nombre total de sommets, indique le degr\'e de connectivit\'e globale du r\'eseau.

Quatri\`emement, le \textbf{diam\`etre} (plus long des plus courts chemins) et la \textbf{longueur moyenne des chemins} sont estim\'es. Le diam\`etre est calcul\'e sur la plus grande composante fortement connexe, car il n'est d\'efini que sur les graphes connexes. Lorsque cette composante d\'epasse 5\,000 sommets, l'excentricit\'e est \'echantillonn\'ee sur 500 n\oe{}uds al\'eatoires (graine fix\'ee \`a 42 pour la reproductibilit\'e) afin de r\'eduire le co\^ut de calcul. La longueur moyenne des chemins est estim\'ee sur la plus grande WCC (version non orient\'ee du graphe) par \'echantillonnage de 1\,000 paires al\'eatoires de sommets.

\subsection{R\'esultats}
\label{sec:rq1_resultats}

\paragraph{M\'etriques de base.}
Le r\'eseau WETH comprend $|V| = 12\,880$ sommets et $|E| = 39\,999$ ar\^etes orient\'ees agr\'eg\'ees (33\,827 ar\^etes non orient\'ees), pour une densit\'e de $d = 2{,}41 \times 10^{-4}$. Le degr\'e moyen est $\bar{k} = 3{,}11$. Ces valeurs t\'emoignent d'un r\'eseau extr\^emement creux~: sur les quelque 165 millions de connexions possibles, seule une infime fraction est effectivement r\'ealis\'ee. La m\'ediane des degr\'es entrants et sortants est de 1{,}0, ce qui signifie que la majorit\'e des adresses n'ont interagi qu'avec un seul autre portefeuille au cours de la journ\'ee.

\paragraph{Distribution des degr\'es.}
La Figure~\ref{fig:rq1_degree} pr\'esente la CCDF des degr\'es entrants et sortants en \'echelle logarithmique. Les deux distributions pr\'esentent une queue lourde caract\'eristique~: le degr\'e entrant maximal atteint $k_{\text{in}}^{\max} = 865$ tandis que le degr\'e sortant maximal culmine \`a $k_{\text{out}}^{\max} = 2\,003$. La d\'ecroissance approximativement lin\'eaire de la CCDF en \'echelle log-log sugg\`ere un comportement compatible avec une distribution \`a queue lourde, typique des r\'eseaux sans \'echelle. Toutefois, nous ne qualifions pas formellement cette distribution de \og loi de puissance \fg{} en l'absence de test statistique d\'edi\'e (test de Kolmogorov-Smirnov, estimation du param\`etre $\alpha$ et de $x_{\min}$)~\cite{clauset2009powerlaw}. L'asym\'etrie entre degr\'es entrants et sortants, avec un maximum sortant plus de deux fois sup\'erieur au maximum entrant, indique la pr\'esence de n\oe{}uds diffuseurs (routeurs DEX, ponts inter-cha\^ines) dont l'activit\'e d'\'emission d\'epasse celle de r\'eception.

\begin{figure}[H]
  \centering
  \includegraphics[width=0.85\textwidth]{rq1_degree_distribution.png}
  \caption{Distribution compl\'ementaire cumulative (CCDF) des degr\'es entrants et sortants du r\'eseau WETH sur \'echelle logarithmique. La queue lourde sugg\`ere un comportement h\'et\'erog\`ene typique des r\'eseaux sans \'echelle~\cite{barabasi1999scalefree}.}
  \label{fig:rq1_degree}
\end{figure}

\paragraph{Composantes connexes.}
Le r\'eseau contient 95 composantes faiblement connexes et 7\,632 composantes fortement connexes. La composante g\'eante faiblement connexe regroupe 12\,538 sommets, soit 97{,}3\,\% du r\'eseau~: la quasi-totalit\'e des adresses actives sont reli\'ees par des chemins non orient\'es. En revanche, la composante g\'eante fortement connexe ne contient que 5\,206 sommets (40{,}4\,\%), ce qui r\'ev\`ele une asym\'etrie directionnelle importante~: de nombreuses adresses re\c{c}oivent ou \'emettent des WETH sans r\'eciprocit\'e. Le Tableau~\ref{tab:rq1_components} r\'esume ces r\'esultats. La Figure~\ref{fig:rq1_components} illustre la distribution de la taille des 20 plus grandes composantes pour chaque type, en \'echelle logarithmique en raison de la dominance de la composante g\'eante.

\begin{table}[H]
\centering
\caption{R\'esum\'e des composantes connexes du r\'eseau WETH.}
\label{tab:rq1_components}
\begin{tabular}{lrr}
\toprule
M\'etrique & Faiblement connexes & Fortement connexes \\
\midrule
Nombre de composantes & 95 & 7\,632 \\
Composante g\'eante (sommets) & 12\,538 & 5\,206 \\
Composante g\'eante (\%) & 97{,}3\,\% & 40{,}4\,\% \\
\bottomrule
\end{tabular}
\end{table}

\begin{figure}[H]
  \centering
  \includegraphics[width=0.85\textwidth]{rq1_components.png}
  \caption{Distribution de la taille des 20 plus grandes composantes connexes (faiblement et fortement) du r\'eseau WETH, en \'echelle logarithmique. La composante g\'eante domine nettement dans les deux cas.}
  \label{fig:rq1_components}
\end{figure}

\paragraph{Diam\`etre et longueur des chemins.}
Le diam\`etre du r\'eseau, estim\'e sur la plus grande composante fortement connexe (5\,206 sommets) par \'echantillonnage de 500 excentricit\'es al\'eatoires, est d'environ $D \approx 14$ sauts. La longueur moyenne des chemins, estim\'ee sur la plus grande composante faiblement connexe (12\,538 sommets, graphe non orient\'e) \`a partir de 1\,000 paires al\'eatoires, est de $\bar{\ell} \approx 3{,}97$ sauts. Malgr\'e les 12\,880 sommets du r\'eseau, deux adresses quelconques de la composante g\'eante sont s\'epar\'ees en moyenne par moins de quatre interm\'ediaires, ce qui t\'emoigne d'une propagation efficace des transferts de valeur.

\subsection{Interpr\'etation}
\label{sec:rq1_interpretation}

La densit\'e tr\`es faible ($d = 2{,}41 \times 10^{-4}$) est typique des r\'eseaux de transactions blockchain, o\`u la plupart des portefeuilles n'interagissent qu'avec une ou deux contreparties. Ce r\'eseau creux est pourtant presque enti\`erement connect\'e (WCC \`a 97{,}3\,\%)~: les transferts transitent par un petit nombre de n\oe{}uds centraux (routeurs DEX, ponts inter-cha\^ines, contrats de pr\^et) qui servent de plaques tournantes~\cite{newman2018networks}.

La queue lourde s'accorde avec le m\'ecanisme d'attachement pr\'ef\'erentiel de Barab\'asi et Albert~\cite{barabasi1999scalefree}~: les adresses les plus connect\'ees attirent davantage de nouvelles connexions. Pendant un crash, l'effet s'amplifie~: les protocoles d'\'echange et de liquidation concentrent un nombre croissant de transactions quand les utilisateurs cherchent tous \`a convertir ou transf\'erer leurs WETH en m\^eme temps. L'asym\'etrie entre degr\'e sortant maximal (2\,003) et entrant maximal (865) sugg\`ere que certains contrats (routeurs DEX ou bots de liquidation) ont diffus\'e des WETH vers un tr\`es grand nombre de destinataires.

L'\'ecart entre composante g\'eante faiblement connexe (97{,}3\,\%) et fortement connexe (40{,}4\,\%) correspond \`a une structure \og en n\oe{}ud papillon \fg{} (\textit{bow-tie}), classique sur les grands r\'eseaux orient\'es~: un noyau fortement connect\'e de transferts bidirectionnels, encadr\'e par des branches entrantes (adresses qui re\c{c}oivent sans \'emettre en retour) et des branches sortantes (adresses qui \'emettent sans recevoir). On retrouve cette g\'eom\'etrie sur d'autres r\'eseaux de crypto-actifs.

Enfin, un diam\`etre court ($D \approx 14$) et une longueur moyenne des chemins de $\bar{\ell} \approx 3{,}97$ pour pr\`es de 13\,000 sommets ressemblent d\'ej\`a \`a un comportement petit-monde (test\'e formellement en RQ5)~: n'importe quelle adresse est s\'epar\'ee de n'importe quelle autre par moins de quatre interm\'ediaires, ce qui a pu faciliter la cascade rapide de transferts de panique le 5 ao\^ut.

La Figure~\ref{fig:rq1_hub_network} illustre cette architecture en \'etoile~: le sous-graphe des 20 n\oe{}uds les plus connect\'es et de leurs voisins directs montre la centralisation extr\^eme du r\'eseau autour de quelques hubs (en orange), chacun irradiant vers de nombreuses adresses p\'eriph\'eriques.

\begin{figure}[H]
  \centering
  \includegraphics[width=0.90\textwidth]{rq1_hub_network.png}
  \caption{Sous-graphe des 20 hubs principaux (en orange) du r\'eseau WETH et de leurs voisins directs. La topologie en \'etoile confirme la structure hub-and-spoke caract\'eristique des r\'eseaux sans \'echelle.}
  \label{fig:rq1_hub_network}
\end{figure}

\paragraph{Exploration interactive.}
Le bouton \og RQ1~: Topologie \fg{} du prototype web~\cite{cryptographexplorer_web}
recalcule le degr\'e moyen et maximal, la densit\'e et le nombre de
composantes pour le sous-graphe actuellement affich\'e. Un curseur fait
varier ce sous-graphe de $300$ \`a $12\,880$ n\oe{}uds (les n\oe{}uds de
plus grand degr\'e en premier), et le bouton \og Distribution \fg{} trace
l'histogramme du degr\'e. Le diam\`etre $D \approx 14$ rapport\'e ici reste
la r\'ef\'erence~: le prototype l'estime \`a la vol\'ee sur ce qui est
visible \`a l'\'ecran.

\clearpage
% sections/rq2.tex
% Phase 4: RQ2 Centralite (RQ2-01 a RQ2-06)

\section{RQ2 -- Analyse de la centralit\'e}
\label{sec:rq2}

La distribution des degr\'es \`a queue lourde (RQ1) sugg\`ere que certains n\oe{}uds concentrent un nombre disproportionn\'e de connexions. On cherche ici \`a identifier ces n\oe{}uds par quatre mesures de centralit\'e~\cite{freeman1977centrality}, et \`a d\'eterminer si les m\^emes adresses dominent syst\'ematiquement les quatre classements ou si chaque mesure capture un aspect diff\'erent de l'influence.

\subsection{M\'ethodologie}
\label{sec:rq2_methodo}

Quatre mesures de centralit\'e sont calcul\'ees sur le graphe orient\'e pond\'er\'e $G = (V, E)$ construit \`a partir des $12\,880$ sommets et $39\,999$ ar\^etes agr\'eg\'ees du r\'eseau WETH. Chaque mesure \'eclaire une dimension diff\'erente de l'importance d'un n\oe{}ud~: le nombre de connexions directes, le r\^ole d'interm\'ediaire, la proximit\'e structurelle et l'importance relative dans un mod\`ele de marche al\'eatoire.

\paragraph{Centralit\'e de degr\'e.}
La centralit\'e de degr\'e mesure le nombre de connexions directes d'un sommet, normalis\'e par le nombre maximal de connexions possibles. Pour un graphe orient\'e, on distingue la centralit\'e de degr\'e entrant et sortant~:
\begin{equation}
C_{\text{in}}(v) = \frac{k_{\text{in}}(v)}{|V|-1}, \qquad C_{\text{out}}(v) = \frac{k_{\text{out}}(v)}{|V|-1}
\label{eq:degree_centrality}
\end{equation}
o\`u $k_{\text{in}}(v)$ et $k_{\text{out}}(v)$ d\'esignent respectivement le degr\'e entrant et sortant du sommet $v$, et $|V| = 12\,880$ le nombre total de sommets. Un degr\'e entrant \'elev\'e identifie les adresses recevant des transferts depuis de nombreuses sources (r\'ecepteurs populaires), tandis qu'un degr\'e sortant \'elev\'e signale les adresses \'emettant vers de nombreux destinataires (distributeurs actifs).

\paragraph{Centralit\'e d'interm\'ediarit\'e.}
La centralit\'e d'interm\'ediarit\'e (\textit{betweenness centrality}), introduite par Freeman~\cite{freeman1977centrality}, quantifie le r\^ole d'un sommet comme pont entre les autres paires de sommets~:
\begin{equation}
C_B(v) = \sum_{\substack{s,t \in V \\ s \neq v \neq t}} \frac{\sigma_{st}(v)}{\sigma_{st}}
\label{eq:betweenness}
\end{equation}
o\`u $\sigma_{st}$ est le nombre de plus courts chemins entre $s$ et $t$, et $\sigma_{st}(v)$ le nombre de ces chemins passant par $v$. Un n\oe{}ud \`a forte interm\'ediarit\'e constitue un goulot d'\'etranglement potentiel~: sa suppression allongerait significativement les chemins entre de nombreuses paires de sommets. Le calcul exact de cette mesure est en $O(|V| \cdot |E|)$, ce qui repr\'esente environ $5 \times 10^8$ op\'erations pour notre graphe. Afin de maintenir un temps d'ex\'ecution raisonnable (environ 5 secondes au lieu de plusieurs heures), nous utilisons un \'echantillonnage al\'eatoire de $k = 500$ sommets pivots avec une graine fix\'ee \`a 42 pour assurer la reproductibilit\'e. Cette approximation est standard dans la litt\'erature pour les graphes de plus de $5\,000$ sommets~\cite{newman2018networks}.

\paragraph{Centralit\'e de proximit\'e.}
La centralit\'e de proximit\'e (\textit{closeness centrality}) mesure la proximit\'e moyenne d'un sommet par rapport \`a tous les autres~:
\begin{equation}
C_C(v) = \frac{|V'|-1}{\displaystyle\sum_{u \in V', u \neq v} d(v,u)}
\label{eq:closeness}
\end{equation}
o\`u $d(v,u)$ d\'esigne la longueur du plus court chemin entre $v$ et $u$, et $|V'|$ le nombre de sommets de la composante consid\'er\'ee. Cette mesure n'est d\'efinie que lorsque tous les sommets sont atteignables depuis $v$ (i.e., $d(v,u) < \infty$ pour tout $u$). Or, le r\'eseau WETH pr\'esente 95 composantes faiblement connexes (cf.~section~\ref{sec:rq1_resultats}), rendant le calcul impossible sur le graphe complet. Nous restreignons donc le calcul \`a la plus grande composante faiblement connexe (WCC) en version non orient\'ee, qui contient $12\,538$ sommets, soit $97{,}3\,\%$ du r\'eseau. Un n\oe{}ud \`a forte proximit\'e est structurellement \og au centre \fg{} du r\'eseau, accessible en peu de sauts depuis n'importe quel autre sommet.

\paragraph{PageRank.}
Le PageRank, introduit par Page et al.~\cite{page1999pagerank}, mod\'elise l'importance d'un sommet par une marche al\'eatoire avec redirection. \`A chaque pas, un marcheur al\'eatoire suit une ar\^ete sortante avec probabilit\'e $(1 - \alpha)$ ou se t\'el\'eporte vers un sommet al\'eatoire avec probabilit\'e $\alpha$. Le score PageRank de $v$ correspond \`a la probabilit\'e stationnaire de pr\'esence du marcheur sur $v$. Nous utilisons le facteur d'amortissement standard $\alpha = 0{,}85$ et la variante pond\'er\'ee~: les poids des ar\^etes (volumes WETH agr\'eg\'es) influencent la probabilit\'e de transition, de sorte que les n\oe{}uds recevant des transferts de volume \'elev\'e obtiennent un score sup\'erieur. Contrairement \`a la simple centralit\'e de degr\'e, le PageRank tient compte de la qualit\'e des connexions entrantes~: \^etre connect\'e \`a un n\oe{}ud lui-m\^eme influent a plus d'impact qu'\^etre connect\'e \`a un n\oe{}ud p\'eriph\'erique.

\subsection{R\'esultats}
\label{sec:rq2_resultats}

\paragraph{Centralit\'e de degr\'e.}
L'adresse avec la centralit\'e de degr\'e entrant la plus \'elev\'ee est \texttt{0xe50f...8128c8} ($C_{\text{in}} = 0{,}0672$), qui a re\c{c}u des transferts depuis $865$ adresses distinctes au cours de la journ\'ee. La centralit\'e de degr\'e sortant donne un autre leader~: \texttt{0x2818...cb7116} ($C_{\text{out}} = 0{,}1556$), qui a \'emis vers $2\,003$ adresses distinctes, soit pr\`es du double du maximum entrant. C'est l'effet des n\oe{}uds distributeurs d\'ej\`a rep\'er\'es en RQ1. Sur les dix premi\`eres adresses de chaque classement, sept apparaissent dans les deux top-10 (\texttt{0x45dd...f50608}, \texttt{0x50ea...5cc6d7}, \texttt{0xa4d8...6228d9}, \texttt{0x86f1...e618e0}, \texttt{0xc590...74c02b}, \texttt{0x0000...148dea}, \texttt{0xdef1...6fee57})~: les n\oe{}uds les plus actifs fonctionnent comme des relais bidirectionnels, qui re\c{c}oivent et redistribuent en gros volume.

\paragraph{Centralit\'e d'interm\'ediarit\'e.}
Le classement par interm\'ediarit\'e (estim\'ee avec $k = 500$ pivots) place \texttt{0x2818...cb7116} en t\^ete ($C_B = 0{,}0816$), suivi de \texttt{0xe50f...8128c8} ($C_B = 0{,}0692$) et \texttt{0x86f1...e618e0} ($C_B = 0{,}0606$). Les huit premiers n\oe{}uds par interm\'ediarit\'e figurent \'egalement dans le top-10 du degr\'e (entrant ou sortant), ce qui sugg\`ere que dans ce r\'eseau, les n\oe{}uds les plus connect\'es sont aussi les ponts structurels les plus critiques. L'adresse \texttt{0x1231...6f4eae} se distingue en apparaissant au rang~9 de l'interm\'ediarit\'e sans figurer dans les top-10 de degr\'e, sugg\'erant un r\^ole sp\'ecifique de pont entre communaut\'es malgr\'e un nombre de connexions plus modeste.

\paragraph{Centralit\'e de proximit\'e.}
Calcul\'ee sur la plus grande WCC ($12\,538$ sommets), la centralit\'e de proximit\'e r\'ev\`ele un classement sensiblement diff\'erent. L'adresse \texttt{0x86f1...e618e0} domine avec $C_C = 0{,}4288$, suivie de \texttt{0x45dd...f50608} ($C_C = 0{,}4240$) et \texttt{0xa4d8...6228d9} ($C_C = 0{,}4216$). Ces trois sommets sont accessibles en moyenne en moins de $2{,}4$ sauts depuis n'importe quel autre sommet de la composante g\'eante. En revanche, les deux leaders en degr\'e sortant et interm\'ediarit\'e (\texttt{0x2818...cb7116} et \texttt{0xe50f...8128c8}) ne figurent pas dans le top-10 de proximit\'e, signe qu'un degr\'e \'elev\'e ne garantit pas une position centrale au sens de la distance moyenne. Les n\oe{}uds les plus proches sont ceux connect\'es \`a des voisins eux-m\^emes bien connect\'es, formant un noyau dense au c\oe{}ur du r\'eseau.

\paragraph{PageRank.}
Le classement PageRank pond\'er\'e place \texttt{0xe50f...8128c8} en t\^ete ($PR = 0{,}0425$), ce qui co\"incide avec le leader du degr\'e entrant. Cette convergence est attendue~: le PageRank favorise les n\oe{}uds recevant de nombreux transferts entrants, particuli\`erement lorsque ceux-ci proviennent d'adresses elles-m\^emes influentes. L'adresse \texttt{0xa4d8...6228d9} occupe la deuxi\`eme position ($PR = 0{,}0280$), suivie de \texttt{0x50ea...5cc6d7} ($PR = 0{,}0226$) et \texttt{0x45dd...f50608} ($PR = 0{,}0223$). Parmi les dix premiers, six figurent \'egalement dans le top-10 du degr\'e entrant, confirmant la corr\'elation mod\'er\'ee entre ces deux mesures ($\rho = 0{,}606$). L'adresse \texttt{0x802b...3f328b} appara\^it au rang~5 du PageRank ($PR = 0{,}0187$) sans figurer dans aucun autre top-10, indiquant qu'elle re\c{c}oit des transferts de volume particuli\`erement \'elev\'e depuis des n\oe{}uds importants.

\paragraph{Analyse comparative.}
Le Tableau~\ref{tab:rq2_deg_pr} pr\'esente le classement comparatif des dix adresses les plus centrales selon la centralit\'e de degr\'e entrant et le PageRank, tandis que le Tableau~\ref{tab:rq2_bet_clo} compare l'interm\'ediarit\'e et la proximit\'e. La Figure~\ref{fig:rq2_barplots} illustre visuellement les top-10 pour les quatre mesures.

\begin{table}[H]
\centering
\caption{Classement comparatif des dix adresses les plus centrales selon la centralit\'e de degr\'e entrant et le PageRank pond\'er\'e.}
\label{tab:rq2_deg_pr}
\begin{tabular}{clclc}
\toprule
Rang & Degr\'e entrant & $C_{\text{in}}$ & PageRank & $PR$ \\
\midrule
1 & \texttt{0xe50f...8128c8} & $0{,}0672$ & \texttt{0xe50f...8128c8} & $0{,}0425$ \\
2 & \texttt{0x86f1...e618e0} & $0{,}0419$ & \texttt{0xa4d8...6228d9} & $0{,}0280$ \\
3 & \texttt{0x45dd...f50608} & $0{,}0415$ & \texttt{0x50ea...5cc6d7} & $0{,}0226$ \\
4 & \texttt{0xc590...74c02b} & $0{,}0366$ & \texttt{0x45dd...f50608} & $0{,}0223$ \\
5 & \texttt{0xa4d8...6228d9} & $0{,}0359$ & \texttt{0x802b...3f328b} & $0{,}0187$ \\
6 & \texttt{0x50ea...5cc6d7} & $0{,}0306$ & \texttt{0xdef1...6fee57} & $0{,}0172$ \\
7 & \texttt{0x0000...148dea} & $0{,}0290$ & \texttt{0x55ca...678207} & $0{,}0160$ \\
8 & \texttt{0xc9e1...35158d} & $0{,}0286$ & \texttt{0x86f1...e618e0} & $0{,}0158$ \\
9 & \texttt{0xdef1...6fee57} & $0{,}0238$ & \texttt{0xb33b...199346} & $0{,}0114$ \\
10 & \texttt{0x9d9d...b5ab47} & $0{,}0232$ & \texttt{0xa6ae...9a9719} & $0{,}0113$ \\
\bottomrule
\end{tabular}
\end{table}

\begin{table}[H]
\centering
\caption{Classement comparatif des dix adresses les plus centrales selon la centralit\'e d'interm\'ediarit\'e (estim\'ee, $k{=}500$) et de proximit\'e (sur la WCC, $12\,538$ sommets).}
\label{tab:rq2_bet_clo}
\begin{tabular}{clclc}
\toprule
Rang & Interm\'ediarit\'e & $C_B$ & Proximit\'e & $C_C$ \\
\midrule
1 & \texttt{0x2818...cb7116} & $0{,}0816$ & \texttt{0x86f1...e618e0} & $0{,}4287$ \\
2 & \texttt{0xe50f...8128c8} & $0{,}0692$ & \texttt{0x45dd...f50608} & $0{,}4240$ \\
3 & \texttt{0x86f1...e618e0} & $0{,}0606$ & \texttt{0xa4d8...6228d9} & $0{,}4216$ \\
4 & \texttt{0x45dd...f50608} & $0{,}0482$ & \texttt{0x50ea...5cc6d7} & $0{,}4192$ \\
5 & \texttt{0xc590...74c02b} & $0{,}0429$ & \texttt{0xbb98...f29be4} & $0{,}4099$ \\
6 & \texttt{0xa4d8...6228d9} & $0{,}0409$ & \texttt{0x4ccd...96f9b1} & $0{,}3883$ \\
7 & \texttt{0x50ea...5cc6d7} & $0{,}0347$ & \texttt{0x95d2...993896} & $0{,}3874$ \\
8 & \texttt{0xec7b...fefda2} & $0{,}0344$ & \texttt{0x3e31...f0ea26} & $0{,}3846$ \\
9 & \texttt{0x1231...6f4eae} & $0{,}0320$ & \texttt{0xe37e...57bd09} & $0{,}3760$ \\
10 & \texttt{0xdef1...6fee57} & $0{,}0311$ & \texttt{0x853e...1f670d} & $0{,}3720$ \\
\bottomrule
\end{tabular}
\end{table}

\begin{figure}[H]
  \centering
  \includegraphics[width=0.95\textwidth]{rq2_centrality_barplots.png}
  \caption{Classement des dix adresses les plus centrales selon chaque mesure de centralit\'e. Les barres horizontales facilitent la comparaison visuelle entre le degr\'e entrant (haut gauche), le degr\'e sortant (haut droit), l'interm\'ediarit\'e (bas gauche) et le PageRank (bas droit).}
  \label{fig:rq2_barplots}
\end{figure}

L'analyse des corr\'elations de Spearman entre les quatre mesures, calcul\'ees sur les $12\,538$ n\oe{}uds communs (intersection avec la WCC), r\'ev\`ele des relations contrast\'ees. La corr\'elation la plus forte s'observe entre le degr\'e entrant et le PageRank ($\rho = 0{,}606$), ce qui est attendu puisque les deux mesures privil\'egient les n\oe{}uds recevant de nombreuses connexions. La corr\'elation entre le degr\'e entrant et l'interm\'ediarit\'e est mod\'er\'ee ($\rho = 0{,}514$)~: les n\oe{}uds les plus connect\'es tendent \`a \^etre des ponts, mais cette relation n'est pas syst\'ematique. En revanche, la corr\'elation entre la proximit\'e et le PageRank est quasi nulle ($\rho = 0{,}008$), ce qui montre que ces deux mesures capturent des aspects tr\`es diff\'erents~: la proximit\'e \'evalue la position g\'eom\'etrique, le PageRank l'importance en termes de flux. La Figure~\ref{fig:rq2_comparison} pr\'esente la matrice compl\`ete de ces corr\'elations.

\begin{figure}[H]
  \centering
  \includegraphics[width=0.75\textwidth]{rq2_centrality_comparison.png}
  \caption{Matrice de corr\'elation de Spearman entre les quatre mesures de centralit\'e, calcul\'ee sur les $12\,538$ n\oe{}uds de la composante g\'eante. Les valeurs \'elev\'ees (rouge fonc\'e) indiquent des classements similaires~; les valeurs proches de z\'ero (blanc) signalent des mesures capturant des dimensions ind\'ependantes de la centralit\'e.}
  \label{fig:rq2_comparison}
\end{figure}

\subsection{Interpr\'etation}
\label{sec:rq2_interpretation}

Cinq \`a sept adresses apparaissent dans les top-10 des quatre mesures. Ces n\oe{}uds \og multi-centraux \fg{} correspondent probablement aux contrats des grands protocoles DEX (QuickSwap, SushiSwap), aux ponts inter-cha\^ines (Polygon Bridge), ou aux protocoles de pr\^et (Aave, Compound). L'adresse \texttt{0x2818...cb7116} est premi\`ere en degr\'e sortant ($C_{\text{out}} = 0{,}1556$, soit 2\,003 destinataires) et en interm\'ediarit\'e ($C_B = 0{,}0816$)~; le profil est celui d'un routeur DEX qui diffuse massivement les WETH tout en servant de pont topologique. \`A l'autre extr\'emit\'e, \texttt{0xe50f...8128c8} domine le degr\'e entrant et le PageRank ($PR = 0{,}0425$), ce qui ressemble \`a un accumulateur de liquidit\'e, probablement un contrat de \textit{staking} ou un \textit{vault} DeFi~\cite{newman2018networks}.

La divergence avec la proximit\'e m\'erite qu'on s'y arr\^ete. Les quatre sommets les plus proches (\texttt{0x86f1...e618e0}, \texttt{0x45dd...f50608}, \texttt{0xa4d8...6228d9}, \texttt{0x50ea...5cc6d7}) forment un noyau dense au c\oe{}ur de la composante g\'eante, accessible en tr\`es peu de sauts. Mais ce noyau ne co\"incide qu'en partie avec les leaders en degr\'e ou en PageRank~: la position g\'eom\'etrique dans le r\'eseau n'est pas la m\^eme dimension que le volume de connexions. La corr\'elation quasi nulle entre proximit\'e et PageRank ($\rho = 0{,}008$) le confirme. Un n\oe{}ud peut \^etre g\'eom\'etriquement central sans concentrer de flux, et inversement.

En chiffres~: environ $8{,}2\,\%$ des plus courts chemins du r\'eseau passent par la seule adresse \texttt{0x2818...cb7116}. Si ce n\oe{}ud avait \'et\'e congestionn\'e ou arr\^et\'e, la propagation des \'echanges aurait \'et\'e nettement ralentie. C\^ot\'e PageRank, \texttt{0xe50f...8128c8} concentre $4{,}25\,\%$ du score total, ce qui correspond \`a son r\^ole d'absorbeur de liquidit\'e pendant le crash.

La Figure~\ref{fig:rq2_network} permet de visualiser ces relations de centralit\'e~: les n\oe{}uds sont dimensionn\'es par leur score PageRank et color\'es selon leur centralit\'e d'interm\'ediarit\'e (du bleu au orange). Les hubs multi-centraux apparaissent comme les plus grands n\oe{}uds aux teintes chaudes, entour\'es de grappes de n\oe{}uds p\'eriph\'eriques.

\begin{figure}[H]
  \centering
  \includegraphics[width=0.90\textwidth]{rq2_centrality_network.png}
  \caption{Sous-graphe des 30 n\oe{}uds les plus centraux et de leurs voisins. La taille des n\oe{}uds refl\`ete le score PageRank et la couleur indique la centralit\'e d'interm\'ediarit\'e (bleu~=~faible, orange~=~\'elev\'ee). Les \'etiquettes identifient les 10 adresses les plus influentes.}
  \label{fig:rq2_network}
\end{figure}

\paragraph{Exploration interactive.}
Le bouton \og RQ2~: Centralit\'e \fg{} du prototype web~\cite{cryptographexplorer_web}
affiche les top-10 du degr\'e et du PageRank lus directement dans le CSV
\texttt{nodes\_metrics.csv}, et relance \`a la demande la \textit{betweenness}
sur un \'echantillon de $50$ pivots (contre $500$ dans ce rapport). Cliquer
sur une adresse surligne ses ar\^etes incidentes, ce qui aide \`a voir si
elle se comporte plut\^ot comme un relais ou comme un accumulateur.

\clearpage
% sections/rq3.tex
% Phase 5: RQ3 Communautes (RQ3-01 a RQ3-06)

\section{RQ3 -- D\'etection de communaut\'es}
\label{sec:rq3}

Le r\'eseau WETH pr\'esente-t-il une structure communautaire, autrement dit des groupes d'adresses qui interagissent pr\'ef\'erentiellement entre elles~? Identifier ces groupes permettrait de retrouver les \'ecosyst\`emes fonctionnels (pools de liquidit\'e, protocoles DeFi, clusters d'utilisateurs) qui structurent les \'echanges sur Polygon. On applique pour cela l'algorithme de Louvain~\cite{blondel2008louvain}, adapt\'e aux r\'eseaux de grande taille.

\subsection{M\'ethodologie}
\label{sec:rq3_methodo}

\paragraph{D\'efinition de la structure communautaire.}
Une communaut\'e dans un r\'eseau est un sous-ensemble de sommets dont la densit\'e d'ar\^etes internes exc\`ede celle attendue dans un graphe al\'eatoire poss\'edant la m\^eme s\'equence de degr\'es. La qualit\'e d'un partitionnement en communaut\'es est mesur\'ee par la \textbf{modularit\'e} $Q$, d\'efinie par~\cite{newman2018networks}~:
\begin{equation}
Q = \frac{1}{2m} \sum_{ij} \left[ A_{ij} - \frac{k_i k_j}{2m} \right] \delta(c_i, c_j)
\label{eq:modularity}
\end{equation}
o\`u $A_{ij}$ est l'\'el\'ement de la matrice d'adjacence, $k_i$ le degr\'e du sommet $i$, $m = |E|$ le nombre total d'ar\^etes, $c_i$ la communaut\'e assign\'ee au sommet $i$, et $\delta$ le symbole de Kronecker ($\delta(c_i, c_j) = 1$ si $c_i = c_j$, $0$ sinon). La modularit\'e compare la fraction d'ar\^etes intracommunautaires \`a celle attendue sous un mod\`ele nul de configuration al\'eatoire. Une valeur de $Q > 0{,}3$ est g\'en\'eralement consid\'er\'ee comme indicatrice d'une structure communautaire significative, tandis que $Q > 0{,}6$ t\'emoigne d'une structure forte.

\paragraph{Algorithme de Louvain.}
L'algorithme de Louvain~\cite{blondel2008louvain} proc\`ede en deux phases it\'eratives. Dans la premi\`ere phase (\textit{optimisation locale}), chaque sommet est d\'eplac\'e vers la communaut\'e voisine qui maximise le gain de modularit\'e $\Delta Q$~; ce processus est r\'ep\'et\'e jusqu'\`a convergence. Dans la seconde phase (\textit{agr\'egation}), les communaut\'es obtenues sont contract\'ees en supern\oe{}uds formant un graphe r\'eduit, et le processus recommence. Ces deux phases alternent jusqu'\`a ce qu'aucun d\'eplacement ne puisse am\'eliorer $Q$. L'algorithme pr\'esente une complexit\'e en $O(|E|)$ par it\'eration, le rendant applicable aux r\'eseaux de grande taille. Nous utilisons l'impl\'ementation native de NetworkX (\texttt{nx.community.louvain\_communities}) avec un param\`etre de r\'esolution $\gamma = 1{,}0$ (r\'esolution standard de Newman-Girvan) et une graine al\'eatoire fix\'ee \`a 42 pour assurer la reproductibilit\'e des r\'esultats.

\paragraph{Conversion en graphe non orient\'e.}
L'algorithme de Louvain op\`ere exclusivement sur des graphes non orient\'es. Le graphe orient\'e pond\'er\'e $G = (V, E)$ construit en RQ1 est donc converti en sa version non orient\'ee~: pour chaque paire de sommets $(u, v)$ connect\'ee par une ou plusieurs ar\^etes orient\'ees, une ar\^ete non orient\'ee est cr\'e\'ee dont le poids est la somme des poids des ar\^etes orient\'ees correspondantes. Cette transformation pr\'eserve la structure de connectivit\'e et les volumes d'\'echange, tout en \'eliminant l'information directionnelle. Ce choix est standard dans la litt\'erature pour appliquer les m\'ethodes de d\'etection de communaut\'es fond\'ees sur la modularit\'e aux r\'eseaux orient\'es~\cite{newman2018networks}. Le graphe non orient\'e r\'esultant comporte $12\,880$ sommets et $33\,827$ ar\^etes pond\'er\'ees.

\paragraph{Strat\'egie de visualisation.}
Le r\'eseau complet de $12\,880$ sommets produirait une visualisation illisible. Nous adoptons donc une strat\'egie de sous-\'echantillonnage cibl\'e~: les sept plus grandes communaut\'es sont s\'electionn\'ees, et pour chacune, un \'echantillon al\'eatoire de 100 n\oe{}uds au maximum est extrait (graine fix\'ee \`a 42), pour un total d'environ 400 n\oe{}uds repr\'esentatifs. Le sous-graphe induit par ces n\oe{}uds est trac\'e avec un algorithme de positionnement \`a ressorts (Fruchterman-Reingold), les n\oe{}uds \'etant color\'es par communaut\'e d'appartenance et dimensionn\'es proportionnellement \`a leur degr\'e dans le sous-graphe.

\subsection{R\'esultats}
\label{sec:rq3_resultats}

\paragraph{Communaut\'es d\'etect\'ees.}
L'algorithme de Louvain identifie 174 communaut\'es dans le r\'eseau WETH, avec une modularit\'e $Q = 0{,}6445$. C'est au-dessus du seuil usuel de $0{,}3$ et proche de $0{,}7$, ce qui indique une structure communautaire bien marqu\'ee~: les adresses forment des groupes au sein desquels les \'echanges sont beaucoup plus fr\'equents que dans un graphe al\'eatoire de m\^eme s\'equence de degr\'es. Parmi ces 174 communaut\'es, 26 contiennent au moins 10 n\oe{}uds, tandis que seulement 3 sont des singletons (communaut\'es r\'eduites \`a une seule adresse). La taille moyenne des communaut\'es est de $74{,}0$ n\oe{}uds, mais la m\'ediane n'est que de $2{,}0$, r\'ev\'elant une distribution des tailles extr\^emement asym\'etrique.

\paragraph{Distribution des tailles.}
La distribution des tailles de communaut\'es pr\'esente une h\'et\'erog\'en\'eit\'e marqu\'ee~: la plus grande communaut\'e regroupe $2\,981$ n\oe{}uds, soit $23{,}1\,\%$ de l'ensemble du r\'eseau, tandis que la majorit\'e des communaut\'es ne contiennent que quelques adresses. Les cinq plus grandes communaut\'es concentrent \`a elles seules $82{,}5\,\%$ des n\oe{}uds du r\'eseau, ce qui t\'emoigne d'une organisation fortement hi\'erarchis\'ee o\`u quelques macro-communaut\'es coexistent avec de nombreux micro-groupes p\'eriph\'eriques. Le Tableau~\ref{tab:rq3_sizes} d\'etaille la composition des dix plus grandes communaut\'es. La Figure~\ref{fig:rq3_sizes} illustre la distribution compl\`ete des vingt plus grandes communaut\'es en \'echelle logarithmique, mettant en \'evidence le rapport de taille sup\'erieur \`a $100\times$ entre la plus grande et la vingti\`eme communaut\'e.

\begin{table}[H]
\centering
\caption{Taille des dix plus grandes communaut\'es d\'etect\'ees par l'algorithme de Louvain ($Q = 0{,}6445$, r\'esolution $\gamma = 1{,}0$).}
\label{tab:rq3_sizes}
\begin{tabular}{crc}
\toprule
Communaut\'e & N\oe{}uds & \% du r\'eseau \\
\midrule
C1  & 2\,981 & 23{,}1\,\% \\
C2  & 2\,414 & 18{,}7\,\% \\
C3  & 2\,338 & 18{,}2\,\% \\
C4  & 1\,542 & 12{,}0\,\% \\
C5  & 1\,381 & 10{,}7\,\% \\
C6  & 805   & 6{,}3\,\% \\
C7  & 179   & 1{,}4\,\% \\
C8  & 176   & 1{,}4\,\% \\
C9  & 138   & 1{,}1\,\% \\
C10 & 101   & 0{,}8\,\% \\
\bottomrule
\end{tabular}
\end{table}

\begin{figure}[H]
  \centering
  \includegraphics[width=0.85\textwidth]{rq3_community_sizes.png}
  \caption{Distribution de la taille des vingt plus grandes communaut\'es, ordonn\'ees par taille d\'ecroissante (\'echelle logarithmique). Les cinq premi\`eres communaut\'es (barres color\'ees) repr\'esentent $82{,}5\,\%$ du r\'eseau~; les communaut\'es restantes (gris) forment une longue tra\^ine de micro-groupes. La ligne pointill\'ee indique la taille moyenne ($74{,}0$ n\oe{}uds).}
  \label{fig:rq3_sizes}
\end{figure}

\paragraph{Structure communautaire.}
La Figure~\ref{fig:rq3_communities} pr\'esente la visualisation du sous-graphe des sept principales communaut\'es (environ 400 n\oe{}uds sous-\'echantillonn\'es). Les communaut\'es apparaissent comme des amas distincts et relativement bien s\'epar\'es dans l'espace de positionnement \`a ressorts, confirmant visuellement la valeur \'elev\'ee de la modularit\'e. Les ar\^etes intercommunautaires, bien que pr\'esentes, se concentrent sur un nombre restreint de n\oe{}uds de forte connectivit\'e situ\'es \`a la p\'eriph\'erie de chaque amas, jouant le r\^ole de ponts entre communaut\'es. Les n\oe{}uds de degr\'e \'elev\'e (repr\'esent\'es par des cercles de grande taille) tendent \`a occuper des positions centrales au sein de leur communaut\'e, coh\'erent avec la structure en \'etoile attendue autour des contrats de protocoles DeFi.

\begin{figure}[H]
  \centering
  \includegraphics[width=0.85\textwidth]{rq3_communities.png}
  \caption{Visualisation du sous-graphe des sept principales communaut\'es d\'etect\'ees par l'algorithme de Louvain~\cite{blondel2008louvain}. Les n\oe{}uds sont color\'es par communaut\'e d'appartenance et dimensionn\'es selon leur degr\'e. Le positionnement \`a ressorts (Fruchterman-Reingold, graine~42) r\'ev\`ele une s\'eparation nette entre les amas communautaires.}
  \label{fig:rq3_communities}
\end{figure}

\paragraph{Concentration et in\'egalit\'e.}
L'\'ecart entre la taille moyenne ($74{,}0$) et la taille m\'ediane ($2{,}0$) des communaut\'es traduit une in\'egalit\'e structurelle extr\^eme~: la distribution des tailles est domin\'ee par quelques communaut\'es g\'eantes, tandis que la majorit\'e des 174 groupes sont des micro-communaut\'es de deux \`a dix adresses. Ce profil est coh\'erent avec la distribution des degr\'es \`a queue lourde identifi\'ee en RQ1~: les n\oe{}uds concentrateurs (routeurs DEX, ponts, protocoles de pr\^et) attirent un grand nombre de voisins dans leur communaut\'e d'appartenance, cr\'eant des amas de grande taille autour de chaque infrastructure DeFi majeure.

\subsection{Interpr\'etation}
\label{sec:rq3_interpretation}

Avec $Q = 0{,}6445$, le r\'eseau WETH s'organise en compartiments fonctionnels plut\^ot qu'en espace d'\'echange homog\`ene. Chaque grande communaut\'e correspond probablement \`a un \'ecosyst\`eme protocolaire (pool de liquidit\'e DEX, protocole de pr\^et, pont inter-cha\^ines, agr\'egateur de rendement). Les adresses d'un m\^eme protocole partagent beaucoup d'ar\^etes, d'o\`u la densit\'e intracommunautaire \'elev\'ee. Le fait que cinq communaut\'es regroupent $82{,}5\,\%$ du r\'eseau indique que l'activit\'e WETH du 5 ao\^ut s'est concentr\'ee sur cinq infrastructures DeFi majeures, les 169 autres communaut\'es regroupant des usages plus marginaux.

Cette structure modulaire change la fa\c{c}on dont un choc se propage. \`A l'int\'erieur d'une communaut\'e, la densit\'e \'elev\'ee fait que la panique se diffuse vite~: les contreparties d'un m\^eme pool se passent le mouvement de vente en quelques sauts. Entre communaut\'es, la propagation passe presque exclusivement par les n\oe{}uds \`a forte interm\'ediarit\'e identifi\'es en RQ2 (\texttt{0x2818...cb7116}, \texttt{0xe50f...8128c8})~; ces ponts intercommunautaires sont moins nombreux, donc le choc met plus de temps \`a franchir les fronti\`eres protocolaires. La compartimentalisation joue donc le r\^ole d'un amortisseur partiel.

La structure en n\oe{}ud papillon de RQ1 se superpose au partitionnement~: le noyau fortement connexe (SCC, $40{,}4\,\%$) se r\'epartit dans les cinq grandes communaut\'es, et les branches entrantes ou sortantes du \textit{bow-tie} correspondent aux micro-communaut\'es p\'eriph\'eriques. Les n\oe{}uds multi-centraux de RQ2 occupent les positions de liaison entre communaut\'es, ce qui explique pourquoi leur interm\'ediarit\'e est \'elev\'ee.

\paragraph{Exploration interactive.}
Le prototype web~\cite{cryptographexplorer_web} reprend directement la
partition Louvain calcul\'ee ici. La colonne \texttt{community} du CSV
embarqu\'e donne une couleur \`a chaque n\oe{}ud, et on rep\`ere alors
d'un coup d'\oe{}il les sept plus grandes communaut\'es et les quelques
adresses qui font le pont entre elles. Le bouton \og RQ3 \fg{} rappelle
$Q = 0{,}6445$ et le nombre de communaut\'es d\'etect\'ees.

\clearpage
% sections/rq4.tex
% Phase 6: RQ4 Temporel (RQ4-01 a RQ4-06)

\section{RQ4 -- Dynamique temporelle}
\label{sec:rq4}

Jusqu'ici, le r\'eseau a \'et\'e analys\'e comme un instantan\'e agr\'eg\'e sur la journ\'ee enti\`ere. Mais le crash du 5 ao\^ut ne s'est pas produit d'un seul coup~: les march\'es asiatiques, europ\'eens et am\'ericains ont r\'eagi les uns apr\`es les autres, cr\'eant des vagues d'activit\'e d\'ecal\'ees. On d\'esagr\`ege ici les donn\'ees par fen\^etre horaire pour voir comment le r\'eseau a \'evolu\'e au fil de la journ\'ee~: quand l'activit\'e a-t-elle atteint son pic, et comment cette signature temporelle se relie-t-elle aux sessions des march\'es financiers mondiaux~\cite{newman2018networks}~?

\subsection{M\'ethodologie}
\label{sec:rq4_methodo}

\paragraph{D\'ecoupage en fen\^etres horaires.}
Le jeu de donn\'ees couvre la journ\'ee compl\`ete du 5 ao\^ut 2024 en temps universel coordonn\'e (\texttt{UTC}). Chaque transaction porte un horodatage Unix (colonne \texttt{chrono}, exprim\'e en secondes depuis l'\'epoque), converti en datetime \texttt{UTC}-aware via \texttt{pd.to\_datetime(unit='s', utc=True)}. Le choix explicite de \texttt{UTC} \'evite toute contamination par le fuseau horaire local de la machine d'ex\'ecution, garantissant la reproductibilit\'e des r\'esultats ind\'ependamment de la localisation g\'eographique de l'analyste. La journ\'ee est partitionn\'ee en 24 fen\^etres horaires $W_h$ ($h \in \{0, 1, \ldots, 23\}$), o\`u $W_h$ regroupe les transactions dont l'horodatage tombe dans l'intervalle $[h{:}00{:}00,\; h{:}59{:}59]$ \texttt{UTC}.

\paragraph{Construction de graphes par fen\^etre.}
Pour chaque heure $h$, un graphe orient\'e pond\'er\'e $G_h = (V_h, E_h)$ est construit \`a partir des transactions tombant dans $W_h$. Les ar\^etes multiples entre une m\^eme paire d'adresses $(u, v)$ sont agr\'eg\'ees par somme des poids WETH, produisant une ar\^ete unique pond\'er\'ee par le volume cumul\'e. Ce proc\'ed\'e g\'en\`ere 24 instantan\'es ind\'ependants du r\'eseau, chacun capturant la structure de connectivit\'e et les flux de valeur propres \`a son heure.

\paragraph{M\'etriques calcul\'ees par fen\^etre.}
Six m\'etriques sont calcul\'ees pour chaque graphe horaire $G_h$~:
\begin{itemize}
  \item \textbf{Nombre de n\oe{}uds actifs}~: $|V_h|$ (adresses distinctes ayant particip\'e \`a au moins un transfert pendant l'heure $h$).
  \item \textbf{Nombre d'ar\^etes}~: $|E_h|$ (paires ordonn\'ees distinctes $(u, v)$ apr\`es agr\'egation).
  \item \textbf{Nombre de transactions}~: $T_h$ (nombre brut de transferts individuels avant agr\'egation).
  \item \textbf{Densit\'e}~: $\delta_h = \frac{|E_h|}{|V_h| \cdot (|V_h| - 1)}$ pour un graphe orient\'e~\cite{newman2018networks}. Cette m\'etrique mesure la fraction de connexions r\'ealis\'ees par rapport au maximum th\'eorique.
  \item \textbf{Degr\'e moyen}~: $\bar{k}_h = \frac{2\,|E_h|}{|V_h|}$, rapportant le nombre d'ar\^etes au nombre de n\oe{}uds.
  \item \textbf{Volume WETH total}~: $\mathcal{V}_h = \sum_{e \in E_h} w(e)$, la somme des poids de toutes les ar\^etes de $G_h$.
\end{itemize}

\paragraph{Gestion des heures vides et identification des pics.}
Toute heure sans transaction produit $G_h = \emptyset$ avec l'ensemble des m\'etriques fix\'ees \`a z\'ero, assurant la robustesse du calcul. L'identification de la \textbf{p\'eriode de pic} repose sur un seuil de 50\,\% du nombre maximal de transactions~: les heures cons\'ecutives dont $T_h$ exc\`ede ce seuil d\'efinissent la plage de pic soutenu. Ce crit\`ere simple distingue les pics isol\'es d'une activit\'e durablement \'elev\'ee.

\subsection{R\'esultats}
\label{sec:rq4_resultats}

\paragraph{Couverture temporelle.}
Le jeu de donn\'ees contient $401\,709$ transactions r\'eparties sur la totalit\'e des 24 heures du 5 ao\^ut 2024 (\texttt{UTC}), de 00h00m05s \`a 23h59m59s. Aucune heure n'est vide~: les 24 fen\^etres pr\'esentent une activit\'e non nulle, t\'emoignant du caract\`ere continu du march\'e des crypto-actifs. Le volume total \'echang\'e sur la journ\'ee s'\'el\`eve \`a $210\,979{,}67$ WETH.

\paragraph{Heures de pic.}
L'activit\'e suit un double pic. Le pic de transactions tombe \`a 13h~UTC avec $28\,146$ transactions et $1\,960$ n\oe{}uds actifs~; le pic de volume tombe \`a 01h~UTC avec $21\,517{,}99$~WETH \'echang\'es et $5\,791$ ar\^etes (maximum de la journ\'ee). Pic transactionnel et pic volum\'etrique ne co\"incident pas, ce qui pointe vers deux r\'egimes d'activit\'e distincts au cours du crash. Le Tableau~\ref{tab:rq4_peaks} pr\'esente les cinq heures les plus actives en nombre de transactions.

\begin{table}[H]
\centering
\caption{Les cinq heures les plus actives du r\'eseau WETH le 5 ao\^ut 2024 (\texttt{UTC}), class\'ees par nombre de transactions.}
\label{tab:rq4_peaks}
\begin{tabular}{crrrr}
\toprule
Heure (\texttt{UTC}) & Transactions & N\oe{}uds & Ar\^etes & Volume WETH \\
\midrule
13h & $28\,146$ & $1\,960$ & $4\,692$ & $16\,699{,}53$ \\
14h & $23\,455$ & $1\,893$ & $4\,568$ & $8\,663{,}91$ \\
01h & $22\,905$ & $1\,781$ & $5\,791$ & $21\,517{,}99$ \\
16h & $20\,056$ & $1\,695$ & $4\,198$ & $4\,603{,}21$ \\
00h & $19\,993$ & $1\,532$ & $4\,974$ & $14\,000{,}83$ \\
\bottomrule
\end{tabular}
\end{table}

La p\'eriode de pic soutenu, d\'efinie par les heures cons\'ecutives d\'epassant 50\,\% du maximum transactionnel ($14\,073$ transactions), s'\'etend de 10h \`a 18h~UTC, soit une plage de 9 heures continues avec une moyenne de $19\,866$ transactions par heure et un volume cumul\'e de $74\,961$~WETH. Les heures nocturnes (00h--07h~UTC) montrent aussi une activit\'e soutenue, qui d\'epasse souvent le seuil de 50\,\% mais sans se rattacher \`a la plage diurne.

\paragraph{\'Evolution des m\'etriques structurelles.}
La Figure~\ref{fig:rq4_temporal} pr\'esente l'\'evolution des six m\'etriques sur 24 heures. Le sous-graphique sup\'erieur montre l'activit\'e du r\'eseau~: les transactions (barres) et les n\oe{}uds/ar\^etes (lignes) suivent un profil en U invers\'e avec un creux relatif \`a 09h \texttt{UTC} ($10\,691$ transactions) encadr\'e par deux p\'eriodes d'activit\'e \'elev\'ee. Le nombre de n\oe{}uds actifs culmine \`a 13h ($1\,960$), co\"incidant avec le pic transactionnel, tandis que le nombre d'ar\^etes atteint son maximum \`a 01h ($5\,791$), co\"incidant avec le pic volum\'etrique.

Le sous-graphique inf\'erieur montre un effet contre-intuitif~: la densit\'e \'evolue en sens inverse du nombre de transactions. L'heure la plus dense est 23h ($\delta = 0{,}002905$), heure de faible activit\'e ($9\,790$ transactions)~; l'heure la moins dense est 13h ($\delta = 0{,}001222$), c'est-\`a-dire le pic transactionnel. L'explication tient \`a l'afflux de nouveaux participants. \`A 13h, $1\,960$ n\oe{}uds distincts sont actifs contre $969$ \`a 23h, mais ces nouveaux arrivants cr\'eent surtout des connexions uniques (un seul transfert)~; $|V_h|$ cro\^it alors plus vite que $|E_h|$, ce qui dilue la densit\'e. Le degr\'e moyen confirme cette tendance~: il atteint son maximum \`a 05h ($\bar{k} = 6{,}52$) et son minimum \`a 15h ($\bar{k} = 4{,}56$), indiquant que les heures nocturnes pr\'esentent un r\'eseau plus compact o\`u chaque n\oe{}ud \'echange avec davantage de contreparties.

\begin{figure}[H]
  \centering
  \includegraphics[width=0.85\textwidth]{rq4_temporal_evolution.png}
  \caption{\'Evolution temporelle de l'activit\'e du r\'eseau WETH sur 24 heures (5 ao\^ut 2024, \texttt{UTC}). En haut~: nombre de transactions (barres) et n\oe{}uds/ar\^etes actifs (lignes). En bas~: densit\'e, degr\'e moyen et volume WETH \'echang\'e. La zone ombr\'ee rouge indique la p\'eriode de pic soutenu (10h--18h \texttt{UTC}).}
  \label{fig:rq4_temporal}
\end{figure}

\paragraph{Carte de chaleur.}
La Figure~\ref{fig:rq4_heatmap} pr\'esente la carte de chaleur normalis\'ee (min-max par m\'etrique) des six m\'etriques sur 24 heures. La repr\'esentation rend lisibles les corr\'elations et les d\'ecalages entre m\'etriques. Transactions, n\oe{}uds et ar\^etes suivent des profils proches, avec des valeurs \'elev\'ees entre 00h et 06h puis entre 10h et 18h. Le volume WETH se distingue par une concentration marqu\'ee \`a 01h (valeur normalis\'ee de 1{,}0), bien au-dessus des heures voisines. La densit\'e suit un gradient invers\'e~: ses valeurs les plus \'elev\'ees (couleurs chaudes) tombent aux heures de faible activit\'e transactionnelle (22h--23h), ce qui rejoint l'observation quantitative pr\'ec\'edente. Le degr\'e moyen tient un profil interm\'ediaire, avec des valeurs \'elev\'ees r\'eparties sur les heures nocturnes (00h--07h) et une diminution progressive pendant le pic diurne.

\begin{figure}[H]
  \centering
  \includegraphics[width=0.85\textwidth]{rq4_heatmap.png}
  \caption{Carte de chaleur de l'activit\'e du r\'eseau WETH par heure et par m\'etrique. Les valeurs sont normalis\'ees entre 0 et 1 pour chaque m\'etrique (min-max). Les heures de forte activit\'e transactionnelle (couleurs chaudes sur les premi\`eres lignes) co\"incident avec une densit\'e et un degr\'e moyen plus faibles, r\'ev\'elant l'effet de dilution par les nouveaux participants.}
  \label{fig:rq4_heatmap}
\end{figure}

\subsection{Interpr\'etation}
\label{sec:rq4_interpretation}

Le double pic (volume \`a 01h~UTC, transactions \`a 13h~UTC) s'inscrit dans la chronologie du Lundi Noir. Le d\'enouement des positions de \textit{carry trade} sur le yen a d\'ebut\'e lors de la s\'eance asiatique (Tokyo ouvre \`a 00h~UTC en heure d'\'et\'e). Les premi\`eres liquidations automatiques de positions \`a effet de levier ont produit des transferts de volume unitaire \'elev\'e, ce qui explique le pic de volume \`a 01h ($21\,517{,}99$~WETH) et le nombre maximal d'ar\^etes ($5\,791$) avec moins de n\oe{}uds qu'au pic diurne. Ces transferts ressemblent \`a des liquidations de protocoles DeFi et \`a des r\'e\'equilibrages de pools, op\'er\'es par un nombre restreint d'adresses de contrats intelligents qui d\'eplacent de gros volumes. Le pic transactionnel de 13h tombe lorsque les march\'es europ\'eens sont en s\'eance et que Wall Street ouvre (13h30~UTC). Il correspond \`a une r\'eaction de masse~: un grand nombre d'adresses individuelles ($1\,960$ n\oe{}uds, maximum de la journ\'ee) font des transferts de volume moyen plus faible, ce qu'on attend d'une panique retail. Le creux relatif \`a 09h ($10\,691$ transactions) co\"incide avec la fermeture de la session asiatique et l'ouverture progressive de la session europ\'eenne, un intervalle transitoire typique des march\'es financiers mondiaux.

L'inversion densit\'e/activit\'e \'eclaire la r\'eponse du r\'eseau au choc. Aux heures de pic, le nombre de n\oe{}uds actifs cro\^it plus vite que celui des ar\^etes~: les nouveaux participants, qu'ils soient attir\'es par la volatilit\'e ou pouss\'es par des liquidations, ouvrent surtout des connexions uniques avec un petit nombre de contreparties (typiquement un routeur DEX ou un pool de liquidit\'e). Cette injection de n\oe{}uds p\'eriph\'eriques dilue la densit\'e du graphe horaire. La structure en n\oe{}ud papillon de RQ1 (section~\ref{sec:rq1_interpretation}) explique pourquoi~: les branches entrantes et sortantes du \textit{bow-tie} s'\'etendent pendant le pic, et y ajoutent des adresses \`a usage unique qui ne participent pas aux \'echanges r\'eciproques du noyau fortement connexe. Aux heures nocturnes (22h--23h), le r\'eseau redevient plus compact et plus dense, domin\'e par les op\'erateurs habituels (les n\oe{}uds multi-centraux de RQ2, section~\ref{sec:rq2}) qui gardent des connexions multiples entre eux.

Le double pic s'inscrit dans la lecture des autres analyses. Les n\oe{}uds \`a forte interm\'ediarit\'e de RQ2, en particulier \texttt{0x2818...cb7116}, ont vraisemblablement canalis\'e les flux de panique entre les communaut\'es d\'etect\'ees en RQ3. La structure communautaire a jou\'e un r\^ole de barri\`ere partielle~: le choc initi\'e dans les protocoles de liquidation nocturne (pic \`a 01h) n'a atteint les traders individuels qu'\`a l'ouverture des sessions europ\'eenne et am\'ericaine (pic \`a 13h).

\paragraph{Exploration interactive.}
RQ4 ne tourne pas vraiment dans le prototype web~\cite{cryptographexplorer_web}~:
le CSV embarqu\'e \texttt{nodes\_metrics.csv} ne contient que des
m\'etriques agr\'eg\'ees par n\oe{}ud, sans \textit{timestamp}. Le bouton
\og RQ4 \fg{} se contente donc de rappeler le double pic 01h/13h en
annotation. Pour rendre cette analyse interactive, il faudrait embarquer
dans le navigateur le CSV brut des $401\,709$ transactions~; c'est l'une
des pistes list\'ees en perspectives (section~\ref{subsec:perspectives}).

\clearpage
% sections/rq5.tex
% Phase 7: RQ5 Petit-Monde (RQ5-01 a RQ5-06)

\section{RQ5 -- Propri\'et\'es petit-monde}
\label{sec:rq5}

Le r\'eseau WETH pr\'esente-t-il les propri\'et\'es d'un \textbf{petit monde}~? Le mod\`ele de Watts et Strogatz~\cite{watts1998smallworld} d\'efinit un tel r\'eseau par un coefficient de clustering bien sup\'erieur \`a celui d'un graphe al\'eatoire de m\^emes param\`etres, combin\'e \`a une longueur moyenne des chemins comparable. On le quantifie par le quotient $\sigma = (C / C_{\text{rand}}) \,/\, (L / L_{\text{rand}})$~\cite{humphries2008smallworld}~: un $\sigma \gg 1$ confirme le caract\`ere petit-monde. Si c'est le cas, les chocs de valeur peuvent se propager rapidement \`a travers le r\'eseau tout en s'amplifiant localement dans des voisinages denses --- un m\'ecanisme favorable aux cascades de liquidations.

\subsection{M\'ethodologie}
\label{sec:rq5_methodo}

\paragraph{Pr\'eparation du graphe.}
L'analyse petit-monde op\`ere sur la version non orient\'ee du graphe de transferts WETH, conform\'ement \`a la d\'efinition originale de Watts et Strogatz~\cite{watts1998smallworld} qui s'applique aux r\'eseaux non orient\'es. Le graphe orient\'e pond\'er\'e $G = (V, E)$ construit en RQ1 est converti en graphe non orient\'e par fusion des ar\^etes r\'eciproques, produisant un graphe de $n = 12\,880$ sommets et $33\,827$ ar\^etes non orient\'ees. La densit\'e du graphe non orient\'e est $p = 0{,}000408$. Ce graphe n'est pas enti\`erement connexe~: les calculs de longueur de chemin sont donc effectu\'es sur la plus grande composante connexe (LCC), qui regroupe $12\,538$ n\oe{}uds ($97{,}3\,\%$ du graphe total) et $33\,576$ ar\^etes, assurant la repr\'esentativit\'e des r\'esultats.

\paragraph{Coefficients de clustering.}
Deux mesures de clustering sont calcul\'ees, capturant des aspects diff\'erents de la densit\'e locale~:
\begin{itemize}
  \item \textbf{Transitivit\'e globale}~: $C_{\text{trans}} = \frac{3 \times |\text{triangles}|}{|\text{triplets}|}$, rapport entre le nombre de triangles ferm\'es et le nombre total de triplets dans le graphe. Cette m\'etrique globale pond\`ere implicitement chaque triangle par sa contribution \`a l'ensemble du r\'eseau.
  \item \textbf{Clustering moyen}~: $\bar{C} = \frac{1}{n} \sum_{v \in V} C_v$, o\`u $C_v$ est la fraction de paires de voisins du sommet $v$ qui sont connect\'es entre eux. Cette mesure locale pond\`ere \'egalement chaque n\oe{}ud, ind\'ependamment de son degr\'e.
\end{itemize}
Ces deux m\'etriques peuvent diverger significativement sur les r\'eseaux \`a distribution de degr\'es h\'et\'erog\`ene~\cite{newman2018networks}~: les n\oe{}uds de tr\`es fort degr\'e tendent \`a avoir un clustering local faible (il est peu probable que tous leurs voisins soient connect\'es entre eux), ce qui abaisse la transitivit\'e globale par rapport au clustering moyen. Pour le calcul du quotient $\sigma$, la \textbf{transitivit\'e} est utilis\'ee de mani\`ere coh\'erente pour le graphe observ\'e et pour les graphes al\'eatoires de r\'ef\'erence, garantissant une comparaison homog\`ene.

\paragraph{Estimation de la longueur des chemins.}
La longueur moyenne des plus courts chemins est estim\'ee sur la LCC ($12\,538$ n\oe{}uds) par \'echantillonnage de $1\,000$ paires al\'eatoires de sommets (graine fix\'ee \`a 42 pour la reproductibilit\'e). Pour chaque paire $(u, v)$, le plus court chemin est calcul\'e via l'algorithme de Dijkstra non pond\'er\'e. Cette approche par \'echantillonnage fournit une estimation fiable sans le co\^ut prohibitif du calcul exhaustif de toutes les paires ($\approx 78$ millions de chemins)~\cite{newman2018networks}. La totalit\'e des $1\,000$ paires \'echantillonn\'ees ont produit un chemin valide, confirmant la pleine connexit\'e de la LCC.

\paragraph{Comparaison avec le mod\`ele Erd\H{o}s-R\'enyi.}
Pour \'etablir une r\'ef\'erence al\'eatoire, 10 graphes sont g\'en\'er\'es selon le mod\`ele $G(n, p)$ d'Erd\H{o}s et R\'enyi~\cite{erdos1959random}, avec les m\^emes param\`etres que le graphe observ\'e~: $n = 12\,880$ sommets et $p = 0{,}000408$ (probabilit\'e de connexion \'egale \`a la densit\'e du graphe observ\'e). Chaque graphe al\'eatoire est g\'en\'er\'e avec une graine d\'eterministe ($42 + i$ pour le $i$-\`eme graphe) assurant la reproductibilit\'e. Pour chaque graphe $G_{\text{ER}}^{(i)}$, la transitivit\'e et la longueur moyenne des chemins (estim\'ee par \'echantillonnage de $1\,000$ paires sur la plus grande composante connexe) sont calcul\'ees. Les valeurs de r\'ef\'erence $C_{\text{rand}}$ et $L_{\text{rand}}$ sont obtenues par moyenne arithm\'etique sur les 10 r\'ealisations.

\paragraph{Quotient sigma.}
Le quotient petit-monde $\sigma$ est calcul\'e manuellement selon la formule de Humphries et Gurney~\cite{humphries2008smallworld}~:
\begin{equation}
\sigma = \frac{C_{\text{obs}} \,/\, C_{\text{rand}}}{L_{\text{obs}} \,/\, L_{\text{rand}}}
\label{eq:sigma}
\end{equation}
o\`u $C_{\text{obs}}$ et $L_{\text{obs}}$ d\'esignent respectivement la transitivit\'e et la longueur moyenne des chemins du graphe observ\'e, et $C_{\text{rand}}$, $L_{\text{rand}}$ celles du mod\`ele al\'eatoire. Un $\sigma \gg 1$ confirme les propri\'et\'es petit-monde~: le clustering est nettement sup\'erieur au hasard tandis que les chemins restent comparables. Ce calcul est effectu\'e manuellement plut\^ot que via la fonction \texttt{nx.sigma()} de NetworkX, dont l'impl\'ementation g\'en\`ere 100 graphes al\'eatoires et s'av\`ere prohibitive en temps de calcul pour un graphe de cette taille.

\subsection{R\'esultats}
\label{sec:rq5_resultats}

\paragraph{Coefficients de clustering.}
La transitivit\'e globale du r\'eseau WETH vaut $C_{\text{trans}} = 0{,}029715$ et le clustering moyen $\bar{C} = 0{,}114295$. L'\'ecart d'un facteur $3{,}8$ entre les deux mesures vient de l'h\'et\'erog\'en\'eit\'e des degr\'es~: les n\oe{}uds de faible degr\'e, majoritaires dans le r\'eseau, ont souvent un clustering local \'elev\'e parce que leurs quelques voisins forment des cliques compl\`etes, ce qui tire $\bar{C}$ au-dessus de la transitivit\'e globale. La distribution des coefficients par n\oe{}ud (Figure~\ref{fig:rq5_clustering_dist}) est bimodale~: $10\,159$ n\oe{}uds ($78{,}9\,\%$) ont un clustering nul, principalement des n\oe{}uds pendants (degr\'e 1) pour lesquels le coefficient n'est pas d\'efini ou est fix\'e \`a z\'ero. Les $2\,721$ n\oe{}uds restants ($21{,}1\,\%$) ont un clustering strictement positif, et une fraction non n\'egligeable d'entre eux approche $C_v = 1{,}0$ (petites cliques int\'egralement interconnect\'ees).

\paragraph{Longueur des chemins.}
La longueur moyenne des chemins, estim\'ee sur la LCC de $12\,538$ n\oe{}uds \`a partir de $1\,000$ paires al\'eatoires, est de $L = 3{,}9700$ sauts. Pour un r\'eseau de pr\`es de $13\,000$ sommets, c'est peu, et c'est conforme au diam\`etre court ($D \approx 14$) d\'ej\`a rapport\'e en RQ1.

\paragraph{Comparaison avec le mod\`ele al\'eatoire.}
Le Tableau~\ref{tab:rq5_comparison} synth\'etise la comparaison entre le r\'eseau observ\'e et la moyenne des 10 graphes Erd\H{o}s-R\'enyi de m\^emes param\`etres. La Figure~\ref{fig:rq5_comparison} illustre visuellement cette comparaison.

\begin{table}[H]
\centering
\caption{Comparaison des propri\'et\'es petit-monde entre le r\'eseau WETH observ\'e et la moyenne de 10 graphes al\'eatoires Erd\H{o}s-R\'enyi de m\^emes param\`etres ($n = 12\,880$, $p = 0{,}000408$).}
\label{tab:rq5_comparison}
\begin{tabular}{lrrr}
\toprule
M\'etrique & R\'eseau observ\'e & ER al\'eatoire (moy.) & Ratio \\
\midrule
Transitivit\'e ($C$) & $0{,}029715$ & $0{,}000381$ & $77{,}94\times$ \\
Longueur moy. chemins ($L$) & $3{,}9700$ & $5{,}9034$ & $0{,}67\times$ \\
\midrule
$\sigma = (C/C_{\text{rand}}) \,/\, (L/L_{\text{rand}})$ & \multicolumn{3}{c}{$\mathbf{115{,}90}$} \\
\bottomrule
\end{tabular}
\end{table}

La transitivit\'e du r\'eseau observ\'e ($C = 0{,}029715$) vaut 78 fois celle des graphes al\'eatoires ($C_{\text{rand}} = 0{,}000381$)~: les adresses ferment beaucoup de triangles, autrement dit les partenaires d'\'echange d'une m\^eme adresse \'echangent souvent entre eux. La longueur moyenne des chemins ($L = 3{,}97$) est 33\,\% plus courte que dans le mod\`ele al\'eatoire ($L_{\text{rand}} = 5{,}90$), parce que les n\oe{}uds centraux servent de raccourcis.

Le quotient $\sigma = 115{,}90$ d\'epasse largement le seuil $\sigma = 1$~: le r\'eseau WETH pr\'esente des propri\'et\'es petit-monde marqu\'ees.

\begin{figure}[H]
  \centering
  \includegraphics[width=0.85\textwidth]{rq5_small_world_comparison.png}
  \caption{Comparaison des propri\'et\'es petit-monde entre le r\'eseau WETH observ\'e (bleu) et la moyenne de 10 graphes Erd\H{o}s-R\'enyi (violet). \`A gauche~: coefficient de clustering (transitivit\'e). \`A droite~: longueur moyenne des chemins. Le r\'eseau observ\'e pr\'esente un clustering $78\times$ sup\'erieur au mod\`ele al\'eatoire avec des chemins $33\,\%$ plus courts, confirmant un quotient $\sigma = 115{,}90$.}
  \label{fig:rq5_comparison}
\end{figure}

\begin{figure}[H]
  \centering
  \includegraphics[width=0.85\textwidth]{rq5_clustering_distribution.png}
  \caption{Distribution des coefficients de clustering locaux par n\oe{}ud. Les lignes verticales indiquent la transitivit\'e globale ($C_{\text{trans}} = 0{,}0297$, rouge) et le clustering moyen ($\bar{C} = 0{,}1143$, jaune). La majorit\'e des n\oe{}uds ($78{,}9\,\%$) pr\'esentent un clustering nul (n\oe{}uds p\'eriph\'eriques \`a connexion unique), tandis qu'une fraction significative approche $C_v = 1$ (cliques locales compl\`etes).}
  \label{fig:rq5_clustering_dist}
\end{figure}

\subsection{Interpr\'etation}
\label{sec:rq5_interpretation}

$\sigma = 115{,}90$ place le r\'eseau WETH tr\`es au-del\`a du seuil petit-monde~\cite{watts1998smallworld}. Le clustering est sup\'erieur de deux ordres de grandeur \`a celui d'un graphe al\'eatoire comparable, donc les adresses transigent souvent avec les partenaires de leurs partenaires. Les chemins restent courts, et m\^eme plus courts que dans le mod\`ele al\'eatoire, parce que les n\oe{}uds centraux (routeurs DEX, ponts, protocoles de pr\^et) servent de raccourcis.

\`A l'\'echelle d'un choc, cette g\'eom\'etrie a deux cons\'equences. D'abord, un mouvement de panique se propage tr\`es vite localement~: quand une adresse liquide, ses contreparties (qui sont elles-m\^emes connect\'ees entre elles) subissent presque imm\'ediatement la m\^eme pression. Le voisinage fonctionne alors comme une chambre d'\'echo. Ensuite, les chemins courts permettent au choc local de gagner le reste du r\'eseau via les ponts intercommunautaires de RQ2. C'est le m\^eme m\'ecanisme \`a deux vitesses que celui d\'ecrit en RQ3, vu sous un autre angle~: diffusion rapide \`a l'int\'erieur d'une communaut\'e, transmission entre communaut\'es en quelques sauts via les hubs.

Les autres analyses convergent dans ce sens. Le noyau fortement connexe \`a $40{,}4\,\%$ de RQ1 et le diam\`etre court fournissent le substrat topologique du petit-monde. Les n\oe{}uds \`a forte interm\'ediarit\'e de RQ2 maintiennent les chemins courts malgr\'e la compartimentalisation observ\'ee en RQ3 ($Q = 0{,}6445$). Et le double pic temporel de RQ4 montre concr\`etement comment un choc initi\'e dans une fen\^etre horaire (les liquidations nocturnes) atteint l'ensemble du r\'eseau quelques heures plus tard.

La Figure~\ref{fig:rq5_network} rend visible cette structure petit-monde~: les n\oe{}uds sont color\'es par leur coefficient de clustering local, du bleu (faible) au rouge (\'elev\'e). Les zones denses \`a forte coloration chaude correspondent aux voisinages coh\'esifs o\`u se forment les triangles locaux, tandis que les n\oe{}uds bleus agissent comme des ponts entre ces grappes.

\begin{figure}[H]
  \centering
  \includegraphics[width=0.90\textwidth]{rq5_clustering_network.png}
  \caption{\'Echantillon du r\'eseau WETH color\'e par coefficient de clustering local (bleu~=~faible, rouge~=~\'elev\'e). Les grappes \`a fort clustering (en rouge) illustrent les voisinages coh\'esifs caract\'eristiques des r\'eseaux petit-monde.}
  \label{fig:rq5_network}
\end{figure}

\paragraph{Exploration interactive.}
Le bouton \og RQ5~: Small-World \fg{} du prototype web~\cite{cryptographexplorer_web}
calcule le clustering et la longueur moyenne des chemins sur le sous-graphe
affich\'e, g\'en\`ere un graphe al\'eatoire de m\^eme densit\'e pour
comparer, et conclut par un OUI/NON \`a partir de seuils ($C > 0{,}3$ et
$L < 6$). C'est utile pour visualiser le ph\'enom\`ene, mais le $\sigma$
exact ($\sigma = 115{,}90$) rapport\'e ici demande la moyenne sur $10$
tirages Erd\H{o}s-R\'enyi via NetworkX, ce que le navigateur ne fait pas.

\clearpage
% sections/rq6.tex
% Phase 8: RQ6 Flux Ponderes (RQ6-01 a RQ6-06)

\section{RQ6 -- Analyse des flux pond\'er\'es}
\label{sec:rq6}

Les analyses pr\'ec\'edentes traitaient les ar\^etes comme de simples connexions binaires. On place ici leurs poids au centre du regard~: comment les volumes de transfert WETH se distribuent-ils sur le r\'eseau~? Sont-ils r\'epartis de mani\`ere homog\`ene, ou concentr\'es autour de quelques \textit{whales} dont les mouvements pilotent la dynamique \'economique~?

Pour r\'epondre, on reprend deux outils classiques de l'\'economie des in\'egalit\'es, transpos\'es \`a l'analyse des r\'eseaux pond\'er\'es~\cite{newman2018networks}~: le coefficient de Gini, mesure synth\'etique de la concentration d'une distribution, et la courbe de Lorenz, sa repr\'esentation graphique. L'identification des ar\^etes les plus lourdes compl\`ete l'analyse en reliant la concentration statistique \`a des acteurs concrets de l'\'ecosyst\`eme DeFi.

\subsection{M\'ethodologie}
\label{sec:rq6_methodo}

\paragraph{Extraction et filtrage des poids.}
Le graphe orient\'e pond\'er\'e $G = (V, E)$ construit en RQ1 comporte $|E| = 39\,999$ ar\^etes, chacune porteuse d'un attribut \texttt{weight} repr\'esentant le volume total de WETH transf\'er\'e entre la paire d'adresses correspondante (pr\'e-agr\'eg\'e par somme depuis les $401\,709$ transactions brutes). Avant toute analyse distributionnelle, les \textit{transactions poussi\`ere} (\textit{dust transactions}), c'est-\`a-dire les ar\^etes de poids inf\'erieur \`a $10^{-12}$~WETH, sont filtr\'ees. Ce seuil \'elimine les micro-transferts n\'egligeables (pings automatis\'es, approbations \`a valeur nulle, erreurs d'arrondi) qui fausseraient les statistiques descriptives et la mesure de concentration sans porter de signification \'economique. Apr\`es filtrage, $1\,598$ ar\^etes poussi\`ere sont retir\'ees ($4{,}0\,\%$ du total), laissant $n = 38\,401$ ar\^etes pour l'analyse.

\paragraph{Distribution des poids.}
La distribution des poids des ar\^etes est repr\'esent\'ee sur un histogramme \`a \'echelle logarithmique sur l'axe des abscisses, avec 50 intervalles espac\'es logarithmiquement (\texttt{np.logspace}). Le recours \`a l'\'echelle logarithmique est indispensable~: la distribution couvrant plus de dix ordres de grandeur (du minimum $\approx 10^{-12}$ au maximum $\approx 10^{4}$ WETH), un histogramme \`a \'echelle lin\'eaire produirait une seule barre visible~\cite{clauset2009powerlaw}. Les lignes verticales de la moyenne et de la m\'ediane sont superpos\'ees afin de visualiser l'asym\'etrie extr\^eme de la distribution.

\paragraph{Coefficient de Gini.}
Le coefficient de Gini $G$ mesure la concentration d'une distribution de valeurs positives. Soit $w_{(1)} \leq w_{(2)} \leq \ldots \leq w_{(n)}$ les poids des $n$ ar\^etes tri\'es par ordre croissant. Le coefficient est d\'efini par~:
\begin{equation}
G = \frac{2 \displaystyle\sum_{i=1}^{n} i \cdot w_{(i)}}{n \displaystyle\sum_{i=1}^{n} w_{(i)}} - \frac{n+1}{n}
\label{eq:gini}
\end{equation}
Un coefficient $G = 0$ indique une distribution parfaitement \'egalitaire (tous les transferts portent le m\^eme volume), tandis que $G = 1$ signifie une concentration maximale (une seule ar\^ete porte la totalit\'e du volume). Originellement introduit en \'economie pour quantifier les in\'egalit\'es de revenus, le coefficient de Gini s'applique naturellement \`a l'analyse des r\'eseaux pond\'er\'es pour caract\'eriser la concentration des flux de valeur~\cite{newman2018networks}.

\paragraph{Courbe de Lorenz.}
La courbe de Lorenz constitue la repr\'esentation graphique de la concentration mesur\'ee par le coefficient de Gini. Elle trace la \textbf{part cumul\'ee du volume total de WETH} (axe des ordonn\'ees) en fonction de la \textbf{part cumul\'ee des ar\^etes tri\'ees par poids croissant} (axe des abscisses). La diagonale \`a $45$\textdegree{} repr\'esente l'\textbf{\'egalit\'e parfaite}~: si tous les transferts portaient le m\^eme volume, la courbe de Lorenz co\"inciderait avec cette diagonale. L'aire comprise entre la ligne d'\'egalit\'e et la courbe de Lorenz est proportionnelle au coefficient de Gini~: plus cette aire est grande, plus la concentration est forte.

\paragraph{Identification des \textit{whales}.}
Les 10 ar\^etes portant les poids les plus \'elev\'es sont identifi\'ees comme les \textit{whales} du r\'eseau. Pour chacune, l'adresse source, l'adresse cible (tronqu\'ees pour la lisibilit\'e), le volume WETH et le nombre de transactions sous-jacentes sont rapport\'es. La part cumul\'ee du volume total repr\'esent\'ee par ces 10 transferts les plus importants est calcul\'ee, quantifiant ainsi le degr\'e de domination des flux les plus lourds sur l'\'economie du r\'eseau.

\subsection{R\'esultats}
\label{sec:rq6_resultats}

\paragraph{Statistiques descriptives des poids.}
Le volume total \'echang\'e sur les $38\,401$ ar\^etes retenues est de $210\,979{,}67$~WETH. La distribution est tr\`es asym\'etrique~: la moyenne des poids vaut $5{,}4941$~WETH, mais la m\'ediane n'est que de $0{,}0231$~WETH, soit un \'ecart de plus de deux ordres de grandeur. L'\'ecart-type ($107{,}39$~WETH) d\'epasse la moyenne d'un facteur 20, ce qui dit la m\^eme chose autrement. Le poids minimum apr\`es filtrage est de $1{,}73 \times 10^{-12}$~WETH (juste au-dessus du seuil poussi\`ere), le maximum atteint $10\,324{,}46$~WETH. Les centiles donnent la forme de la queue~: $25^{\text{e}}$ \`a $0{,}0018$~WETH, $75^{\text{e}}$ \`a $0{,}2375$~WETH, $90^{\text{e}}$ \`a $1{,}78$~WETH, $95^{\text{e}}$ \`a $5{,}80$~WETH, $99^{\text{e}}$ \`a $65{,}44$~WETH. Autrement dit, 75\,\% des ar\^etes portent moins d'un quart de WETH, alors que le 1\,\% sup\'erieur d\'epasse 65~WETH. La Figure~\ref{fig:rq6_weight_dist} illustre cette distribution en \'echelle logarithmique, o\`u l'\'ecart entre les lignes de moyenne et de m\'ediane mat\'erialise l'asym\'etrie.

\begin{figure}[H]
  \centering
  \includegraphics[width=0.85\textwidth]{rq6_weight_distribution.png}
  \caption{Distribution des poids des transferts WETH (\'echelle logarithmique, 50 intervalles). Les lignes verticales indiquent la moyenne ($5{,}49$ WETH, rouge) et la m\'ediane ($0{,}023$ WETH, jaune). L'\'ecart de plus de deux ordres de grandeur entre ces deux valeurs t\'emoigne de l'extr\^eme asym\'etrie de la distribution.}
  \label{fig:rq6_weight_dist}
\end{figure}

\paragraph{Concentration des flux~: coefficient de Gini.}
Le coefficient de Gini calcul\'e sur les $38\,401$ ar\^etes filtr\'ees s'\'el\`eve \`a $G = 0{,}9780$, une valeur t\'emoignant d'une \textbf{concentration extr\^eme} des flux de valeur. Pour contextualiser, le coefficient de Gini des revenus des pays les plus in\'egalitaires (Afrique du Sud, Br\'esil) se situe entre $0{,}50$ et $0{,}65$~; les pays scandinaves pr\'esentent des Gini de $0{,}25$--$0{,}30$. Le Gini de $0{,}9780$ observ\'e ici d\'epasse largement ces r\'ef\'erences, ce qui est coh\'erent avec la nature des r\'eseaux de transactions blockchain, o\`u les flux sont structurellement domin\'es par des contrats automatiques traitant des volumes consid\'erables.

\paragraph{Courbe de Lorenz.}
La courbe de Lorenz (Figure~\ref{fig:rq6_lorenz}) traduit visuellement la concentration~: elle \'epouse le coin inf\'erieur droit et reste tr\`es loin de la diagonale d'\'egalit\'e. En chiffres, les 90\,\% d'ar\^etes les plus l\'eg\`eres ne portent que 2{,}3\,\% du volume total~; les 10\,\% les plus lourdes en concentrent 97{,}7\,\%. L'asym\'etrie est plus forte encore dans l'extr\^eme sup\'erieur~: le 1\,\% d'ar\^etes les plus lourdes (384 ar\^etes) transporte \`a lui seul 81{,}3\,\% du volume, soit environ $171\,200$~WETH sur $210\,980$.

\begin{figure}[H]
  \centering
  \includegraphics[width=0.85\textwidth]{rq6_lorenz_curve.png}
  \caption{Courbe de Lorenz des flux WETH. La diagonale repr\'esente l'\'egalit\'e parfaite. La courbe, fortement concave et \'epousant le coin inf\'erieur droit, confirme la concentration extr\^eme des volumes ($G = 0{,}9780$). L'aire gris\'ee entre la diagonale et la courbe est proportionnelle au coefficient de Gini.}
  \label{fig:rq6_lorenz}
\end{figure}

\paragraph{Identification des \textit{whales}.}
Le Tableau~\ref{tab:rq6_whales} pr\'esente les 10 transferts les plus volumineux du r\'eseau. Ces 10 ar\^etes p\`esent \`a elles seules 24{,}4\,\% du volume total ($51\,573{,}51$ WETH sur $210\,979{,}67$). En regardant de plus pr\`es, les 10 \textit{whales} forment 5 paires bidirectionnelles. Le transfert \no1 (de \texttt{0x55ca...8207} vers \texttt{0x3cc2...6580}, $10\,324{,}46$~WETH en 252 transactions) et le transfert \no2 (sens inverse, $10\,323{,}64$~WETH en 134 transactions) impliquent les m\^emes adresses dans les deux directions, avec des volumes quasi identiques. Les quatre autres paires suivent le m\^eme sch\'ema. Cette sym\'etrie bidirectionnelle est la signature caract\'eristique des \textbf{pools de liquidit\'e des \'echanges d\'ecentralis\'es} (\textit{AMM}, \textit{Automated Market Makers})~: les contrats de pool re\c{c}oivent du WETH des utilisateurs qui \'echangent, puis en restituent lors des transactions inverses, cr\'eant un flux bidirectionnel de volume \'equivalent.

\begin{table}[H]
\centering
\caption{Top-10 des transferts WETH les plus importants (\textit{whales}). Les rangs impairs et pairs forment des paires bidirectionnelles entre les m\^emes adresses, caract\'eristiques des pools de liquidit\'e DEX.}
\label{tab:rq6_whales}
\begin{tabular}{rlllr}
\toprule
Rang & Source & Cible & Volume (WETH) & Tx \\
\midrule
1 & \texttt{0x55ca...8207} & \texttt{0x3cc2...6580} & $10\,324{,}46$ & 252 \\
2 & \texttt{0x3cc2...6580} & \texttt{0x55ca...8207} & $10\,323{,}64$ & 134 \\
3 & \texttt{0xa6ae...9719} & \texttt{0x3974...214f} & $5\,264{,}12$ & 214 \\
4 & \texttt{0x3974...214f} & \texttt{0xa6ae...9719} & $5\,263{,}40$ & 118 \\
5 & \texttt{0xa4d8...28d9} & \texttt{0x1fd4...2a78} & $3\,805{,}55$ & 213 \\
6 & \texttt{0x1fd4...2a78} & \texttt{0xa4d8...28d9} & $3\,805{,}21$ & 108 \\
7 & \texttt{0xb694...c474} & \texttt{0x3365...4662} & $3\,274{,}71$ & 190 \\
8 & \texttt{0x3365...4662} & \texttt{0xb694...c474} & $3\,274{,}69$ & 107 \\
9 & \texttt{0x9cef...10f2} & \texttt{0x5928...f09f} & $3\,119{,}02$ & 210 \\
10 & \texttt{0x5928...f09f} & \texttt{0x9cef...10f2} & $3\,118{,}72$ & 111 \\
\bottomrule
\end{tabular}
\end{table}

\subsection{Interpr\'etation}
\label{sec:rq6_interpretation}

\paragraph{Concentration extr\^eme et infrastructure DeFi.}
Un Gini de $0{,}9780$ dit qu'une fraction infime des ar\^etes transporte
presque tout le volume. Ce n'est pas un effet de crise propre au 5 ao\^ut~:
c'est la fa\c{c}on dont fonctionne la DeFi en r\'egime normal, o\`u les flux
de valeur transitent par quelques contrats d'infrastructure (pools AMM,
protocoles de pr\^et et de liquidation, ponts inter-cha\^ines) qui traitent
par construction des volumes \'enormes. Les 5 paires bidirectionnelles du
Tableau~\ref{tab:rq6_whales} ressemblent \`a des pools WETH/stablecoin ou
WETH/MATIC sur les principaux DEX de Polygon (QuickSwap, Uniswap~v3,
SushiSwap). La quasi-\'egalit\'e des volumes dans les deux sens ($10\,324{,}46$
contre $10\,323{,}64$~WETH pour la premi\`ere paire) est la signature du
m\'ecanisme AMM~: chaque \'echange dans un sens produit un flux compensatoire
dans l'autre, ce qui maintient la r\'eserve du pool autour de son \'equilibre.

\paragraph{Amplification de crise et cascade de liquidations.}
Le 5 ao\^ut, cette concentration a amplifi\'e le choc plut\^ot que de
l'absorber. Quelques ar\^etes lourdes (les pools de liquidit\'e) ont
canalis\'e l'essentiel des ventes paniqu\'ees. Quand le prix du WETH chute,
les positions \`a effet de levier passent sous leur seuil de liquidation, ce
qui d\'eclenche des transferts automatis\'es vers ces pools, qui propagent
\`a leur tour le choc de prix. La courbe de Lorenz montre l'envers du
d\'ecor~: les 90\,\% d'ar\^etes les plus l\'eg\`eres (petits traders, transferts
secondaires) ne pesent que $2{,}3\,\%$ du volume. L'\'economie du r\'eseau
repose sur l'infrastructure, pas sur le flux retail. \c{C}a colle avec le
double pic de RQ4 (section~\ref{sec:rq4_interpretation})~: le pic de volume
\`a 01h~UTC ($21\,517{,}99$~WETH) tient \`a ces transferts massifs entre pools
(liquidations nocturnes), alors que le pic transactionnel de 13h tient \`a
un grand nombre de petits transferts (panique retail) qui ne d\'eplacent
qu'une fraction modeste du volume.

Les r\'esultats des six RQ pointent dans le m\^eme sens. Le r\'eseau est
\og sans \'echelle \fg{} sur les connexions (RQ1) comme sur les volumes~: les
\textit{whales} de RQ6 sont les m\^emes routeurs et ponts qui dominaient les
classements de centralit\'e en RQ2. La structure modulaire de RQ3 cloisonne
les petits flux mais les transferts massifs franchissent les fronti\`eres
communautaires. Le pic de volume nocturne (RQ4) s'explique par ces ar\^etes
\`a gros d\'ebit, et les chemins courts du petit-monde (RQ5,
$\bar{\ell} \approx 3{,}97$) permettent au choc de liquidit\'e d'atteindre les
r\'egions p\'eriph\'eriques en quelques sauts.

La Figure~\ref{fig:rq6_whale_network} illustre cette concentration~: les 50 ar\^etes les plus lourdes forment un sous-graphe dense o\`u l'\'epaisseur des liens est proportionnelle au volume transf\'er\'e. Les n\oe{}uds sources nettes apparaissent en bleu et les puits nets en rouge.

\begin{figure}[H]
  \centering
  \includegraphics[width=0.90\textwidth]{rq6_whale_network.png}
  \caption{Sous-graphe des 50 transferts WETH les plus volumineux. L'\'epaisseur des ar\^etes est proportionnelle au volume transf\'er\'e (\'echelle logarithmique). Les n\oe{}uds bleus sont des sources nettes (\'emetteurs), les rouges des puits nets (r\'ecepteurs).}
  \label{fig:rq6_whale_network}
\end{figure}

Avec RQ6, l'image d'ensemble se pr\'ecise~: un \'ecosyst\`eme DeFi o\`u l'architecture en \'etoile, les chemins courts et la concentration des volumes autour de quelques pools de liquidit\'e cr\'eent les conditions d'une amplification syst\'emique des chocs financiers.

\paragraph{Exploration interactive.}
Le bouton \og RQ6~: Flux Pond\'er\'es \fg{} du prototype web~\cite{cryptographexplorer_web}
calcule un Gini et la part du top-20\,\% \`a la vol\'ee, mais sur la
distribution des \emph{degr\'es} des n\oe{}uds visibles, pas sur les
volumes WETH (qui ne sont pas dans le CSV embarqu\'e). La vraie valeur
$G = 0{,}9780$, calcul\'ee ici sur les poids des $38\,401$ ar\^etes
filtr\'ees, reste donc la r\'ef\'erence. Pour la reproduire dans le
navigateur, il faudrait embarquer un CSV des poids par ar\^ete et y
brancher le calcul du Gini et de la courbe de Lorenz.

\clearpage
\input{sections/difficultes}
\clearpage
% sections/bilan.tex
% Phase 10: Bilan -- Conclusion et Perspectives (RAPP-07, RAPP-08)

\section{Bilan}
\label{sec:bilan}

\subsection{Conclusion}
\label{subsec:conclusion}

Les 401\,709 transferts WETH du 5 ao\^ut 2024 sur Polygon ont \'et\'e
analys\'es sous six angles compl\'ementaires. Le Tableau~\ref{tab:synthese_rq}
r\'esume les r\'esultats.

\begin{table}[H]
\centering
\caption{Synth\`ese des r\'esultats par question de recherche}
\label{tab:synthese_rq}
\begin{tabular}{llll}
\toprule
\textbf{RQ} & \textbf{M\'etrique cl\'e} & \textbf{R\'esultat}
  & \textbf{Interpr\'etation} \\
\midrule
RQ1 & Densit\'e & $2{,}41 \times 10^{-4}$
  & R\'eseau creux, structure hub \\
RQ2 & Centralit\'e & Top-10 multi-mesures
  & Routeurs DEX dominants \\
RQ3 & Modularit\'e $Q$ & 0{,}6445
  & Compartimentalisation DeFi \\
RQ4 & Pics d'activit\'e & 01h / 13h UTC
  & Liquidations vs panique \\
RQ5 & Sigma $\sigma$ & 115{,}90
  & Petit-monde confirm\'e \\
RQ6 & Gini $G$ & 0{,}9780
  & Concentration extr\^eme \\
\bottomrule
\end{tabular}
\end{table}

Le r\'eseau est creux ($d = 2{,}41 \times 10^{-4}$) mais presque enti\`erement
connect\'e (WCC \`a 97{,}3\,\%). Il s'organise autour de quelques hubs (routeurs
DEX, ponts, coffres-forts de protocoles) qui canalisent l'essentiel des flux
(RQ1, RQ2). Ces hubs structurent 174 communaut\'es ($Q = 0{,}6445$), qui
correspondent aux diff\'erents protocoles DeFi (RQ3). L'analyse temporelle
(RQ4) fait appara\^itre deux pics distincts~: liquidations automatiques \`a
gros volume autour de 01h~UTC, puis panique retail \`a 13h~UTC \`a
l'ouverture de Wall Street. Les m\'etriques petit-monde sortent tr\`es
au-dessus du r\'ef\'erent al\'eatoire ($\sigma = 115{,}90$, RQ5), et la
concentration de volume est extr\^eme ($G = 0{,}9780$, RQ6)~: 1\,\% des
ar\^etes transportent 81{,}3\,\% du volume.

Mis bout \`a bout, ces six r\'esultats dessinent un m\^eme r\'eseau~: rapide et
efficace en temps normal parce que peu de hubs traitent beaucoup de flux,
fragile en temps de crise pour la m\^eme raison. Quelques hubs concentrent
connexions et volumes, les chemins sont courts donc les chocs voyagent vite,
et la liquidit\'e d\'epend d'un petit nombre de pools. Le 5 ao\^ut 2024, ces
propri\'et\'es ont jou\'e ensemble.

\subsection{Perspectives}
\label{subsec:perspectives}

Plusieurs pistes prolongeraient ce travail.

\paragraph{Fitting power-law formel.}
L'analyse topologique (RQ1) a sugg\'er\'e une distribution des degr\'es \`a
queue lourde sans toutefois fournir de test statistique formel. Une analyse
rigoureuse selon la m\'ethodologie de Clauset et
al.~\cite{clauset2009powerlaw} (estimation de l'exposant $\alpha$ et du
seuil $x_{\min}$ par maximum de vraisemblance, suivie d'un test de
Kolmogorov-Smirnov comparant les mod\`eles power-law, lognormal et
exponentiel) permettrait de statuer sur la nature sans \'echelle du
r\'eseau.

\paragraph{Analyse multi-tokens.}
La pr\'esente \'etude se limite aux transferts WETH. \'Etendre l'analyse
aux transferts ETH, USDC et MATIC sur la m\^eme journ\'ee r\'ev\'elerait si
les propri\'et\'es structurelles observ\'ees (topologie hub-and-spoke,
petit-monde, concentration extr\^eme) sont sp\'ecifiques au WETH ou
caract\'eristiques de l'ensemble de l'\'ecosyst\`eme DeFi sur Polygon.

\paragraph{Analyse multi-jours.}
Le clich\'e unique du 5 ao\^ut 2024 capture le crash mais pas la
dynamique de r\'ecup\'eration. Analyser les jours pr\'ec\'edant, pendant et
suivant le Lundi Noir permettrait d'observer l'\'evolution de la topologie
du r\'eseau \`a travers les phases de crise~: stabilit\'e pr\'e-choc,
propagation du choc et restructuration post-choc.

\paragraph{Visualisation interactive.}
On a d\'evelopp\'e en parall\`ele du rapport un prototype web pour explorer
les m\^emes donn\'ees \`a la souris~: \url{https://crypto-graph-explorer.vercel.app}~\cite{cryptographexplorer_web}.
La visualisation est faite en D3.js, l'interface en React. Les six RQ et les
cinq analyses compl\'ementaires de la section~\ref{sec:complementaires} y
sont accessibles depuis un panneau lat\'eral, avec zoom, drag et filtrage
par communaut\'e sur les $12\,880$ sommets. Le prototype tourne dans le
navigateur, donc sur des m\'etriques pr\'e-calcul\'ees et quelques
approximations~: \textit{betweenness} \'echantillonn\'ee, Gini calcul\'e sur
les degr\'es au lieu des volumes WETH, et la dimension temporelle est
absente faute de \textit{timestamps} dans le CSV embarqu\'e. Les chiffres
exacts du rapport ($\sigma$, $Q$, $G$) restent ceux obtenus localement avec
Python. Restent \`a embarquer dans le prototype les \textit{timestamps}
(RQ4) et les poids WETH par ar\^ete (RQ6), et \`a y brancher un
\textit{fitting} formel de la loi de puissance pour la distribution des
degr\'es.


%% Bibliographie et Webographie
\clearpage
\printbibliography[title={Bibliographie}, notkeyword=web]
\clearpage
\printbibliography[title={Webographie}, keyword=web]

%% Annexes
\clearpage
\appendix
% annexes/cahier_charges.tex
% Cahier des charges du projet - Analyse du reseau WETH Polygon

\section{Cahier des charges}
\label{ann:cahier}

Ce cahier des charges d\'efinit le cadre, les objectifs et les contraintes du projet d'analyse du r\'eseau de transferts WETH sur la blockchain Polygon lors du crash du 5 ao\^ut 2024.

% ---- Contexte ----

\subsection{Contexte}

Le 5 ao\^ut 2024, surnomm\'e le \og Lundi Noir \fg{} des march\'es financiers, a \'et\'e marqu\'e par un effondrement g\'en\'eralis\'e des march\'es boursiers et des cryptomonnaies. Ce crash, d\'eclench\'e par le d\'enouement brutal du \textit{carry trade} sur le yen japonais suite \`a la hausse surprise des taux directeurs de la Banque du Japon~\cite{boj2024rate, reuters2024carrytrade}, a provoqu\'e des mouvements massifs de capitaux sur l'ensemble des r\'eseaux blockchain.

Dans ce contexte, la blockchain Polygon --- solution de scalabilit\'e de couche~2 pour Ethereum --- a enregistr\'e une activit\'e exceptionnelle de transferts du token WETH (\textit{Wrapped Ether}). Le pr\'esent projet applique la th\'eorie des graphes \`a l'analyse de ce r\'eseau de transferts, en mod\'elisant chaque adresse comme un n\oe ud et chaque transfert comme une ar\^ete pond\'er\'ee par le montant en WETH.

Le jeu de donn\'ees, extrait de Google BigQuery~\cite{bigquery_public}, contient 401\,709 transactions effectu\'ees sur la journ\'ee du 5 ao\^ut 2024, impliquant 12\,880 adresses uniques et 39\,999 ar\^etes agr\'eg\'ees.

% ---- Objectifs ----

\subsection{Objectifs}

L'objectif principal est de r\'epondre \`a six questions de recherche (RQ) couvrant diff\'erents aspects de l'analyse de r\'eseaux~:

\begin{itemize}
  \item \textbf{RQ1 --- Topologie du r\'eseau}~: analyser la structure globale du graphe (distribution des degr\'es, densit\'e, composantes connexes, diam\`etre) et caract\'eriser sa topologie.

  \item \textbf{RQ2 --- Centralit\'e}~: identifier les n\oe uds les plus influents du r\'eseau \`a l'aide de quatre mesures de centralit\'e (degr\'e, betweenness, closeness, PageRank).

  \item \textbf{RQ3 --- Communaut\'es}~: d\'etecter les communaut\'es au sein du r\'eseau via l'algorithme de Louvain~\cite{blondel2008louvain} et \'evaluer la modularit\'e de la partition obtenue.

  \item \textbf{RQ4 --- Dynamique temporelle}~: \'etudier l'\'evolution du r\'eseau au cours de la journ\'ee en d\'ecoupant les transactions par fen\^etres horaires.

  \item \textbf{RQ5 --- Propri\'et\'es petit-monde}~: v\'erifier si le r\'eseau pr\'esente des propri\'et\'es de type \textit{small-world}~\cite{watts1998smallworld} en comparant ses m\'etriques \`a un graphe al\'eatoire d'Erd\H{o}s--R\'enyi~\cite{erdos1959random}.

  \item \textbf{RQ6 --- Flux pond\'er\'es}~: analyser la distribution des volumes WETH transf\'er\'es, calculer le coefficient de Gini et identifier les flux dominants (\textit{whales}).
\end{itemize}

% ---- Perimetre fonctionnel ----

\subsection{P\'erim\`etre fonctionnel}

\paragraph{Inclus dans le p\'erim\`etre~:}
\begin{itemize}
  \item Analyse compl\`ete des 6 questions de recherche list\'ees ci-dessus
  \item G\'en\'eration de 16 figures PNG haute qualit\'e (200~DPI)
  \item R\'edaction d'un rapport LaTeX complet en fran\c{c}ais, sans limite de pages
  \item Code Python modulaire et reproductible
\end{itemize}

\paragraph{Exclu du p\'erim\`etre~:}
\begin{itemize}
  \item Analyse multi-jours (seule la journ\'ee du 5 ao\^ut 2024 est consid\'er\'ee)
  \item Analyse d'autres tokens que WETH
  \item R\'e-extraction des donn\'ees depuis BigQuery (utilisation exclusive du CSV existant)
\end{itemize}

\paragraph{R\'ealis\'e en d\'epassement du p\'erim\`etre initial~:}
\begin{itemize}
  \item Application web interactive (\textit{prototype} D3.js + React)
    d\'eploy\'ee \`a l'adresse \url{https://crypto-graph-explorer.vercel.app}~\cite{cryptographexplorer_web},
    qui reproduit les six RQ et les cinq analyses compl\'ementaires de la
    section~\ref{sec:complementaires} sur des m\'etriques pr\'e-calcul\'ees.
\end{itemize}

% ---- Contraintes techniques ----

\subsection{Contraintes techniques}

\begin{itemize}
  \item \textbf{Langage}~: Python~3.10+ avec les biblioth\`eques NetworkX~3.4~\cite{networkx_docs}, Matplotlib~3.9, Pandas~2.2, NumPy~2.1, SciPy~1.14
  \item \textbf{Donn\'ees}~: fichier CSV de 74~Mo contenant 401\,709 transactions avec cinq colonnes (\texttt{transaction\_hash}, \texttt{chrono}, \texttt{source}, \texttt{target}, \texttt{weight})
  \item \textbf{Reproductibilit\'e}~: graines al\'eatoires fix\'ees \`a 42 pour tous les algorithmes stochastiques
  \item \textbf{Mod\'elisation}~: graphe orient\'e pond\'er\'e (DiGraph) par d\'efaut, conversion en non orient\'e pour les RQ3 et RQ5
  \item \textbf{Template LaTeX}~: conforme au mod\`ele de l'Universit\'e Paris Nanterre (Licence MIAGE), compilation via pdfLaTeX + Biber
  \item \textbf{Langue du rapport}~: fran\c{c}ais int\'egral avec conventions typographiques respect\'ees via le package \texttt{babel}
\end{itemize}

% ---- Livrables attendus ----

\subsection{Livrables attendus}

\begin{itemize}
  \item \textbf{Rapport PDF}~: document LaTeX complet comprenant introduction, environnement de travail, description du projet, biblioth\`eques utilis\'ees, 6 sections d'analyse (RQ1--RQ6), difficult\'es rencontr\'ees, bilan, bibliographie, webographie et annexes
  \item \textbf{Figures}~: 16 fichiers PNG \`a 200~DPI dans le r\'epertoire \texttt{figures/}, avec une trame graphique uniforme (fonds sombres pour les graphes r\'eseau, fonds clairs pour les graphiques analytiques)
  \item \textbf{Code source Python}~: 8 modules dans \texttt{src/} (6 analyses RQ + constructeur de graphe + module de style) et un orchestrateur \texttt{run\_all.py}
  \item \textbf{Jeu de donn\'ees}~: fichier \texttt{polygon\_weth\_crash\_2024-08-05.csv} (401\,709 transactions, 74~Mo)
\end{itemize}

\clearpage
\input{annexes/execution}
\clearpage
\input{annexes/manuel}

\end{document}
